\section{核验锚点簇：建文、永乐：继承战争、北京与远征（1399--1424）}
\label{register:1399-1424_jianwen_yongle}
继承战争与北方经营共享朝廷文书，却不共享同一条道路。以下锚点按年序收束可核的行动、承担资源的人群和材料的盲区，而不把后来形成的北京秩序倒写为当时的必然。

以下为读者可跳转的首次描写。它们是按王朝年序编排的证据性短条，而非把账本字段伪装成同一种长篇叙事：每项只在所列材料可承受的范围内给出行动者、空间、后果与限度。

\begin{description}
  \item[\eventanchor{EM-1399-JINGNAN-BEGIN}{1399（建文元年）七月：燕王朱棣以“靖难”为名起兵反对建文朝}] 明与燕王集团；朝廷与官员、军队与卫所；跨区域行政、资源或军政网络；不假定同步执行。建文朝先后处理周、齐、湘等王，燕王担忧自身兵权和封国安全 内战使北方卫所、运河和山东交通成为争夺对象；胜负将决定继承叙事的书写者 材料：《明太宗实录》建文元年七月条；《明史·成祖本纪》与《方孝孺传》；北平地方碑刻和方志线索。限度：起兵和主要诏令可靠；双方“奉天靖难”与“削藩”合法性措辞均具政治目的
  \item[\eventanchor{ML-1399-001}{1399（建文元年）六月：建文帝归还朱棣的儿子}] 明；朝廷与官员；跨区域行政、资源或军政网络；不假定同步执行。从1399（建文元年）六月的《惠帝实录》建文元年条；《明史·本纪》及相关列传；诏令、奏疏或会典版本出发，朝廷与官员因“建文帝归还朱棣的儿子”在跨区域的行政、资源与军政网络进入可见的行政程序。 关于“建文帝归还朱棣的儿子”的实际执行与代价，须从1399（建文元年）六月之后的命令和回应另行检验；一份《惠帝实录》建文元年条；《明史·本纪》及相关列传；诏令、奏疏或会典版本的编年叙述并不足以代替它们。 材料：《惠帝实录》建文元年条；《明史·本纪》及相关列传；诏令、奏疏或会典版本。限度：「建文帝归还朱棣的儿子」借《惠帝实录》建文元年条；《明史·本纪》及相关列传；诏令、奏疏或会典版本获得可检索的时间坐标；跨区域行政、资源或军政网络里朝廷与官员的路线、人数与后果若无其他材料相证，均保留为不可见的部分。
  \item[\eventanchor{ML-1399-002}{1399（建文元年）七月：一名军官以煽动叛乱罪逮捕了朱棣的两名下级官员}] 明；朝廷与官员、军队与卫所、流动与反抗集团；跨区域行政、资源或军政网络；不假定同步执行。就「一名军官以煽动叛乱罪逮捕了朱棣的两名下级官员」而言，明与朝廷与官员、军队与卫所、流动与反抗集团在跨区域行政、资源或军政网络的行动并非凭空发生；不过现存材料只标示节点，前序命令、运输和地方响应仍须逐件核验。 「一名军官以煽动叛乱罪逮捕了朱棣的两名下级官员」之后，朝廷与官员、军队与卫所、流动与反抗集团在跨区域行政、资源或军政网络的行动空间有所改写；这种变化须由随后文书和地方材料追踪，不能把初报数字直接当作最终后果。 材料：《惠帝实录》建文元年条；《明史·兵志》及相关列传；边镇、地方志或对方材料。限度：关于「一名军官以煽动叛乱罪逮捕了朱棣的两名下级官员」，《惠帝实录》建文元年条；《明史·兵志》及相关列传；边镇、地方志或对方材料所能确证的是条目所示的时点和行动；跨区域行政、资源或军政网络内朝廷与官员、军队与卫所、流动与反抗集团的规模、意图及社会后果需要异源材料才能判断。
  \item[\eventanchor{ML-1399-003}{1399（建文元年）八月5日：荆南战役：朱棣向邻县发起进攻}] 明；朝廷与官员、军队与卫所、地方社会与精英；地方或省域；未具城镇者不补造路线。「荆南战役：朱棣向邻县发起进攻」将朝廷与官员、军队与卫所、地方社会与精英的选择落在地方或省域；编年记录可显示先后，却不能代替对兵源、道路和行政准备的证明，故不把相邻事件直接写成原因。 「荆南战役：朱棣向邻县发起进攻」使明在地方或省域的下一步选择重新围绕朝廷与官员、军队与卫所、地方社会与精英展开；其影响应落到随后可见的军令、安置或交涉，而非以未经互校的人数概括。 材料：《惠帝实录》建文元年条；《明史·兵志》及相关列传；边镇、地方志或对方材料。限度：就「荆南战役：朱棣向邻县发起进攻」而论，《惠帝实录》建文元年条；《明史·兵志》及相关列传；边镇、地方志或对方材料提供的是1399（建文元年）八月5日的记录线索；朝廷与官员、军队与卫所、地方社会与精英在地方或省域的实际行动细节不能由这一材料链独自补足。
  \item[\eventanchor{ML-1399-004}{1399（建文元年）九月25日：靖难之役：朱棣大败建文帝十三万大军}] 明；朝廷与官员、军队与卫所；跨区域行政、资源或军政网络；不假定同步执行。在跨区域行政、资源或军政网络发生的「靖难之役：朱棣大败建文帝十三万大军」使朝廷与官员、军队与卫所必须作出即时选择；可见记录没有完整交代先前部署，故只把军政背景列为条件，不写成已经证实的原因。 作为后续影响，「靖难之役：朱棣大败建文帝十三万大军」改变了朝廷与官员、军队与卫所在跨区域行政、资源或军政网络继续调度、驻守或交涉的起点；能确认的是条件变化，而不是由单一叙述推算出的完整战果。 材料：《惠帝实录》建文元年条；《明史·兵志》及相关列传；边镇、地方志或对方材料。限度：「靖难之役：朱棣大败建文帝十三万大军」可由《惠帝实录》建文元年条；《明史·兵志》及相关列传；边镇、地方志或对方材料定位到1399（建文元年）九月25日，但这些叙述未完整保存跨区域行政、资源或军政网络的执行与受影响者经验，不能把沉默写成事实。
  \item[\eventanchor{ML-1399-005}{1399（建文元年）十一月12日：靖难之役：建文帝军队围攻北平，三周后被迫撤退}] 明；朝廷与官员、军队与卫所；跨区域行政、资源或军政网络；不假定同步执行。「靖难之役：建文帝军队围攻北平，三周后被迫撤退」所涉朝廷与官员、军队与卫所的行动受跨区域行政、资源或军政网络的空间与资源约束；年序材料保存了节点，却没有保存足以裁定出兵、守备或支援动机的全过程。 就后果说，「靖难之役：建文帝军队围攻北平，三周后被迫撤退」重排了跨区域行政、资源或军政网络中朝廷与官员、军队与卫所可以据守、经过或谈判的条件；现有记载不能据此精确判定兵数、伤亡或地方社会的全部代价。 材料：《惠帝实录》建文元年条；《明史·兵志》及相关列传；边镇、地方志或对方材料。限度：以《惠帝实录》建文元年条；《明史·兵志》及相关列传；边镇、地方志或对方材料核对「靖难之役：建文帝军队围攻北平，三周后被迫撤退」时，只能确认其在1399（建文元年）十一月12日的叙述位置；朝廷与官员、军队与卫所在跨区域行政、资源或军政网络的具体经验仍需由可定位的异类材料补证。
  \item[\eventanchor{EM-1400-BAIGOU}{1400（建文二年）：白沟河战役后燕军突破南下通道}] 明中央军与燕军；朝廷与官员、军队与卫所；跨区域行政、资源或军政网络；不假定同步执行。建文军依赖李景隆等集中兵力围燕；燕军须夺取河北南下道路和补给 燕军获得向山东、直隶推进的机会；中央军统帅和兵源组织的弱点暴露 材料：《明太宗实录》建文二年条；《明史·成祖本纪》与《李景隆传》；北平、保定地方志战场条。限度：战役结局大体可靠；火器作用、军数和将领责任在双方叙事中夸张不一
  \item[\eventanchor{ML-1400-001}{1400（建文二年）一月：靖难之役：朱棣攻入山西}] 明；朝廷与官员、军队与卫所；跨区域行政、资源或军政网络；不假定同步执行。就「靖难之役：朱棣攻入山西」而言，明与朝廷与官员、军队与卫所在跨区域行政、资源或军政网络的行动并非凭空发生；不过现存材料只标示节点，前序命令、运输和地方响应仍须逐件核验。 此后，朝廷与官员、军队与卫所在跨区域行政、资源或军政网络面对的并非原来的行动条件；「靖难之役：朱棣攻入山西」留下需要继续处理的秩序与供给问题，损失和胜负仍应按异源材料分别核对。 材料：《惠帝实录》建文二年条；《明史·兵志》及相关列传；边镇、地方志或对方材料。限度：《惠帝实录》建文二年条；《明史·兵志》及相关列传；边镇、地方志或对方材料只能把「靖难之役：朱棣攻入山西」置入1400（建文二年）一月的年序；对于跨区域行政、资源或军政网络中朝廷与官员、军队与卫所的路线、人数和动机，仍须另据地方或对方材料核验。
  \item[\eventanchor{ML-1400-002}{1400（建文二年）五月18日：靖难之战：朱棣军伤亡惨重}] 明；朝廷与官员、军队与卫所；跨区域行政、资源或军政网络；不假定同步执行。「靖难之战：朱棣军伤亡惨重」出现时，朝廷与官员、军队与卫所是在跨区域行政、资源或军政网络的限制下行动；资料所给的是年序定位而非动机自述，前期军政压力只能作为解释线索，不能充作定论。 「靖难之战：朱棣军伤亡惨重」之后，跨区域行政、资源或军政网络内朝廷与官员、军队与卫所仍须处理由该行动留下的控制、通行或防务问题；本条不把战报中的胜负、斩获或伤亡扩大为无从核对的结论。 材料：《惠帝实录》建文二年条；《明史·兵志》及相关列传；边镇、地方志或对方材料。限度：《惠帝实录》建文二年条；《明史·兵志》及相关列传；边镇、地方志或对方材料使「靖难之战：朱棣军伤亡惨重」的时间位置可追查；至于朝廷与官员、军队与卫所在跨区域行政、资源或军政网络的命令传递、路线和损失，必须保持待核而不能随编年补写。
  \item[\eventanchor{ML-1400-003}{1400（建文二年）六月8日：靖难之役：朱棣围攻德州}] 明；朝廷与官员、军队与卫所；跨区域行政、资源或军政网络；不假定同步执行。发生在跨区域行政、资源或军政网络的「靖难之役：朱棣围攻德州」牵动朝廷与官员、军队与卫所；已知材料让我们看见这一时点，却没有把准备行动、地方资源与决策意图完整连起来，因果只能保留为待核问题。 「靖难之役：朱棣围攻德州」让朝廷与官员、军队与卫所在跨区域行政、资源或军政网络的既有部署出现新的承接点；其中期影响要看后续调度是否落实，不能用胜败标签替代具体结果。 材料：《惠帝实录》建文二年条；《明史·兵志》及相关列传；边镇、地方志或对方材料。限度：对「靖难之役：朱棣围攻德州」的《惠帝实录》建文二年条；《明史·兵志》及相关列传；边镇、地方志或对方材料记录可支持年序核验；它不足以替代跨区域行政、资源或军政网络的地方材料，故不把朝廷与官员、军队与卫所的战报、奏报或叙事当成无争议结果。
  \item[\eventanchor{EM-1401-DONGCHANG}{1401（建文三年）十二月：东昌之战中燕军受挫，张玉战死}] 明中央军与燕军；朝廷与官员、军队与卫所；跨区域行政、资源或军政网络；不假定同步执行。燕军南下后补给线拉长，建文军在山东仍可集结守军反击 燕军一度受阻并调整战略；持续内战增加沿线州县征发和交通损耗 材料：《明太宗实录》建文三年十二月条；《明史·张玉传》；山东地方志东昌战事条。限度：战役与张玉死亡可靠；胜败归因及伤亡数字不宜只用《太宗实录》
  \item[\eventanchor{ML-1401-001}{1401（建文三年）一月9日：靖难之役：朱棣部队在山东遭炸，伤亡惨重，被迫撤退}] 明；朝廷与官员、地方社会与精英；地方或省域；未具城镇者不补造路线。「靖难之役：朱棣部队在山东遭炸，伤亡惨重，被迫撤退」使朝廷与官员、地方社会与精英在地方或省域的处境变得可见；本条未保存从征发到出动的全链条，故以可核条件而非概括性“压力”标示其背景。 发生「靖难之役：朱棣部队在山东遭炸，伤亡惨重，被迫撤退」后，地方或省域仍是朝廷与官员、地方社会与精英必须面对的行动场域；可见结果是问题被带入下一段年序，未见的伤亡、掠夺与居民去留不在本条擅断。 材料：《惠帝实录》建文三年条；《明史·兵志》及相关列传；边镇、地方志或对方材料。限度：《惠帝实录》建文三年条；《明史·兵志》及相关列传；边镇、地方志或对方材料记录「靖难之役：朱棣部队在山东遭炸，伤亡惨重，被迫撤退」的方式限定了可说范围：可追踪1401（建文三年）一月9日的行动节点，不能凭此还原地方或省域中朝廷与官员、地方社会与精英的每一处部署。
  \item[\eventanchor{ML-1401-002}{1401（建文三年）四月5日：靖难之役：朱棣军队在德州附近大败皇军}] 明；朝廷与官员、军队与卫所；跨区域行政、资源或军政网络；不假定同步执行。从跨区域行政、资源或军政网络看「靖难之役：朱棣军队在德州附近大败皇军」，朝廷与官员、军队与卫所的行动连接着当年的军政网络；资料没有充分说明何种先行安排直接促成它，不能把编年相邻当作因果证明。 「靖难之役：朱棣军队在德州附近大败皇军」并未结束朝廷与官员、军队与卫所在跨区域行政、资源或军政网络的压力，而是把问题转到后续防务、运输或行政处置；只有后续记录能说明它实际造成何种改变。 材料：《惠帝实录》建文三年条；《明史·兵志》及相关列传；边镇、地方志或对方材料。限度：《惠帝实录》建文三年条；《明史·兵志》及相关列传；边镇、地方志或对方材料为「靖难之役：朱棣军队在德州附近大败皇军」留下编年入口，却未必保留跨区域行政、资源或军政网络中朝廷与官员、军队与卫所的完整视角；本条因此不据之裁定人数、责任或动机。
  \item[\eventanchor{ML-1401-003}{1401（建文三年）八月：靖难之战：皇军迫使朱棣北撤北平}] 明；朝廷与官员、军队与卫所；跨区域行政、资源或军政网络；不假定同步执行。「靖难之战：皇军迫使朱棣北撤北平」记录了朝廷与官员、军队与卫所在跨区域行政、资源或军政网络的一次处置；其兵员、道路或地方支持并非都见于同一材料，因而不以总括理由覆盖未见环节。 在跨区域行政、资源或军政网络，朝廷与官员、军队与卫所因「靖难之战：皇军迫使朱棣北撤北平」而要重新安排守备、行军或沟通；这一后续压力可被标示，但战果与伤亡不能只凭一方叙事定量。 材料：《惠帝实录》建文三年条；《明史·兵志》及相关列传；边镇、地方志或对方材料。限度：「靖难之战：皇军迫使朱棣北撤北平」借《惠帝实录》建文三年条；《明史·兵志》及相关列传；边镇、地方志或对方材料获得可检索的时间坐标；跨区域行政、资源或军政网络里朝廷与官员、军队与卫所的路线、人数与后果若无其他材料相证，均保留为不可见的部分。
  \item[\eventanchor{ML-1401-004}{1401（建文三年）八月：建文帝限制佛教和道教的土地占有规模}] 明；朝廷与官员、宗教与文化行动者；跨区域行政、资源或军政网络；不假定同步执行。《惠帝实录》建文三年条；《明史·本纪》及相关列传；诏令、奏疏或会典版本的1401（建文三年）八月条并未写成宫廷传记；它所确证的是朝廷与官员、宗教与文化行动者围绕“建文帝限制佛教和道教的土地占有规模”在跨区域的行政、资源与军政网络的行动痕迹。 “建文帝限制佛教和道教的土地占有规模”在1401（建文三年）八月之后的人事、诏令或交涉应另随可见文书追索；《惠帝实录》建文三年条；《明史·本纪》及相关列传；诏令、奏疏或会典版本并不能独自裁定动机、密议和受影响者经验。 材料：《惠帝实录》建文三年条；《明史·本纪》及相关列传；诏令、奏疏或会典版本。限度：「建文帝限制佛教和道教的土地占有规模」虽见于《惠帝实录》建文三年条；《明史·本纪》及相关列传；诏令、奏疏或会典版本，但材料形成者未必观察到跨区域行政、资源或军政网络的全部过程；朝廷与官员、宗教与文化行动者的行动范围和受影响群体不能由此单独决定。
  \item[\eventanchor{ML-1401-005}{1401（建文三年）十月：靖难之战：帝国军队被驱逐出北平地区}] 明；朝廷与官员、军队与卫所；跨区域行政、资源或军政网络；不假定同步执行。围绕「靖难之战：帝国军队被驱逐出北平地区」，明在跨区域行政、资源或军政网络面对的是朝廷与官员、军队与卫所的现场处置；材料确认行动进入编年，却不足以把此前调兵、补给或地方压力串成确定因果链。 「靖难之战：帝国军队被驱逐出北平地区」使明在跨区域行政、资源或军政网络的下一步选择重新围绕朝廷与官员、军队与卫所展开；其影响应落到随后可见的军令、安置或交涉，而非以未经互校的人数概括。 材料：《惠帝实录》建文三年条；《明史·兵志》及相关列传；边镇、地方志或对方材料。限度：就「靖难之战：帝国军队被驱逐出北平地区」而论，《惠帝实录》建文三年条；《明史·兵志》及相关列传；边镇、地方志或对方材料提供的是1401（建文三年）十月的记录线索；朝廷与官员、军队与卫所在跨区域行政、资源或军政网络的实际行动细节不能由这一材料链独自补足。
  \item[\eventanchor{EM-1402-NANJING-FALL}{1402（建文四年）六月：燕军入南京，建文帝去向不明，朱棣即位}] 明与燕军；朝廷与官员、军队与卫所；跨区域行政、资源或军政网络；不假定同步执行。燕军控制长江北岸与运河通道，建文朝内部指挥和防御体系瓦解 永乐朝开始重编前朝记录并重置人事；建文遗臣与继承合法性问题长期存在 材料：《明太宗实录》建文四年六月条；《明史·成祖本纪》《惠帝本纪》；南京明故宫考古与地方传说材料。限度：入城和即位可靠；宫城火灾、建文帝生死及出逃传说不能裁作确定事实
  \item[\eventanchor{ML-1402-001}{1402（建文四年）一月：靖难之战：朱棣攻克山东西北}] 明；朝廷与官员、军队与卫所、地方社会与精英；地方或省域；未具城镇者不补造路线。从地方或省域看「靖难之战：朱棣攻克山东西北」，朝廷与官员、军队与卫所、地方社会与精英的行动连接着当年的军政网络；资料没有充分说明何种先行安排直接促成它，不能把编年相邻当作因果证明。 「靖难之战：朱棣攻克山东西北」让朝廷与官员、军队与卫所、地方社会与精英在地方或省域的既有部署出现新的承接点；其中期影响要看后续调度是否落实，不能用胜败标签替代具体结果。 材料：《惠帝实录》建文四年条；《明史·兵志》及相关列传；边镇、地方志或对方材料。限度：对「靖难之战：朱棣攻克山东西北」的《惠帝实录》建文四年条；《明史·兵志》及相关列传；边镇、地方志或对方材料记录可支持年序核验；它不足以替代地方或省域的地方材料，故不把朝廷与官员、军队与卫所、地方社会与精英的战报、奏报或叙事当成无争议结果。
  \item[\eventanchor{ML-1402-002}{1402（建文四年）三月3日：靖难之役：朱棣攻克徐州}] 明；朝廷与官员、军队与卫所；跨区域行政、资源或军政网络；不假定同步执行。「靖难之役：朱棣攻克徐州」使朝廷与官员、军队与卫所在跨区域行政、资源或军政网络的处境变得可见；本条未保存从征发到出动的全链条，故以可核条件而非概括性“压力”标示其背景。 对明而言，「靖难之役：朱棣攻克徐州」使跨区域行政、资源或军政网络的资源和安全议题进入下一轮处置；本条所能承受的是这一转折，不是对战场损失或地方反应的完整裁断。 材料：《惠帝实录》建文四年条；《明史·兵志》及相关列传；边镇、地方志或对方材料。限度：「靖难之役：朱棣攻克徐州」在1402（建文四年）三月3日的出现可回查《惠帝实录》建文四年条；《明史·兵志》及相关列传；边镇、地方志或对方材料；它们不等于跨区域行政、资源或军政网络的全景记录，朝廷与官员、军队与卫所未被记下的选择与损失不作推定。
  \item[\eventanchor{ML-1402-003}{1402（建文四年）四月：靖难之役：朱棣在苏州大败皇军}] 明；朝廷与官员、军队与卫所；跨区域行政、资源或军政网络；不假定同步执行。对于明的「靖难之役：朱棣在苏州大败皇军」，跨区域行政、资源或军政网络与朝廷与官员、军队与卫所限定了行动的可行范围；史料未逐项记录前置命令和供给，不能以后来结果补造一个整齐的起因。 此后，朝廷与官员、军队与卫所在跨区域行政、资源或军政网络面对的并非原来的行动条件；「靖难之役：朱棣在苏州大败皇军」留下需要继续处理的秩序与供给问题，损失和胜负仍应按异源材料分别核对。 材料：《惠帝实录》建文四年条；《明史·兵志》及相关列传；边镇、地方志或对方材料。限度：《惠帝实录》建文四年条；《明史·兵志》及相关列传；边镇、地方志或对方材料只能把「靖难之役：朱棣在苏州大败皇军」置入1402（建文四年）四月的年序；对于跨区域行政、资源或军政网络中朝廷与官员、军队与卫所的路线、人数和动机，仍须另据地方或对方材料核验。
  \item[\eventanchor{ML-1402-004}{1402（建文四年）五月23日：靖难之役：朱棣在安徽被皇军击退}] 明；朝廷与官员、军队与卫所；跨区域行政、资源或军政网络；不假定同步执行。「靖难之役：朱棣在安徽被皇军击退」所涉朝廷与官员、军队与卫所的行动受跨区域行政、资源或军政网络的空间与资源约束；年序材料保存了节点，却没有保存足以裁定出兵、守备或支援动机的全过程。 就后果说，「靖难之役：朱棣在安徽被皇军击退」重排了跨区域行政、资源或军政网络中朝廷与官员、军队与卫所可以据守、经过或谈判的条件；现有记载不能据此精确判定兵数、伤亡或地方社会的全部代价。 材料：《惠帝实录》建文四年条；《明史·兵志》及相关列传；边镇、地方志或对方材料。限度：以《惠帝实录》建文四年条；《明史·兵志》及相关列传；边镇、地方志或对方材料核对「靖难之役：朱棣在安徽被皇军击退」时，只能确认其在1402（建文四年）五月23日的叙述位置；朝廷与官员、军队与卫所在跨区域行政、资源或军政网络的具体经验仍需由可定位的异类材料补证。
  \item[\eventanchor{ML-1402-005}{1402（建文四年）五月28日：靖难之役：朱棣在灵璧大败皇军}] 明；朝廷与官员、军队与卫所；跨区域行政、资源或军政网络；不假定同步执行。在跨区域行政、资源或军政网络发生的「靖难之役：朱棣在灵璧大败皇军」使朝廷与官员、军队与卫所必须作出即时选择；可见记录没有完整交代先前部署，故只把军政背景列为条件，不写成已经证实的原因。 作为后续影响，「靖难之役：朱棣在灵璧大败皇军」改变了朝廷与官员、军队与卫所在跨区域行政、资源或军政网络继续调度、驻守或交涉的起点；能确认的是条件变化，而不是由单一叙述推算出的完整战果。 材料：《惠帝实录》建文四年条；《明史·兵志》及相关列传；边镇、地方志或对方材料。限度：「靖难之役：朱棣在灵璧大败皇军」可由《惠帝实录》建文四年条；《明史·兵志》及相关列传；边镇、地方志或对方材料定位到1402（建文四年）五月28日，但这些叙述未完整保存跨区域行政、资源或军政网络的执行与受影响者经验，不能把沉默写成事实。
  \item[\eventanchor{ML-1402-006}{1402（建文四年）六月7日：靖难之役：朱棣大军渡过淮河}] 明；朝廷与官员、军队与卫所；跨区域行政、资源或军政网络；不假定同步执行。就1402（建文四年）六月7日而言，“靖难之役：朱棣大军渡过淮河”让朝廷与官员、军队与卫所在跨区域的行政、资源与军政网络的资源安排进入《惠帝实录》建文四年条；《明会典》相关条目；地方志、黄册/契约/河工材料的可见范围。 “靖难之役：朱棣大军渡过淮河”只能把1402（建文四年）六月7日的线索接到地方志、契约、河工或仓储材料；《惠帝实录》建文四年条；《明会典》相关条目；地方志、黄册/契约/河工材料所载的中央条文不能直接界定执行范围。 材料：《惠帝实录》建文四年条；《明会典》相关条目；地方志、黄册/契约/河工材料。限度：《惠帝实录》建文四年条；《明会典》相关条目；地方志、黄册/契约/河工材料记录「靖难之役：朱棣大军渡过淮河」的方式限定了可说范围：可追踪1402（建文四年）六月7日的行动节点，不能凭此还原跨区域行政、资源或军政网络中朝廷与官员、军队与卫所的每一处部署。
  \item[\eventanchor{ML-1402-007}{1402（建文四年）六月17日：靖难之役：朱棣拿下扬州}] 明；朝廷与官员；跨区域行政、资源或军政网络；不假定同步执行。「靖难之役：朱棣拿下扬州」把朝廷与官员带到跨区域行政、资源或军政网络的具体关口；现存年条只留下这一行动节点，前期兵力、粮道和地方协力的次序未逐项保存，不能倒推为单一动因。 「靖难之役：朱棣拿下扬州」之后，朝廷与官员在跨区域行政、资源或军政网络的行动空间有所改写；这种变化须由随后文书和地方材料追踪，不能把初报数字直接当作最终后果。 材料：《惠帝实录》建文四年条；《明史·兵志》及相关列传；边镇、地方志或对方材料。限度：关于「靖难之役：朱棣拿下扬州」，《惠帝实录》建文四年条；《明史·兵志》及相关列传；边镇、地方志或对方材料所能确证的是条目所示的时点和行动；跨区域行政、资源或军政网络内朝廷与官员的规模、意图及社会后果需要异源材料才能判断。
  \item[\eventanchor{ML-1402-008}{1402（建文四年）七月1日：靖难之役：朱棣被拦在南京对面的长江边}] 明；朝廷与官员；跨区域行政、资源或军政网络；不假定同步执行。「靖难之役：朱棣被拦在南京对面的长江边」出现时，朝廷与官员是在跨区域行政、资源或军政网络的限制下行动；资料所给的是年序定位而非动机自述，前期军政压力只能作为解释线索，不能充作定论。 在跨区域行政、资源或军政网络，朝廷与官员因「靖难之役：朱棣被拦在南京对面的长江边」而要重新安排守备、行军或沟通；这一后续压力可被标示，但战果与伤亡不能只凭一方叙事定量。 材料：《惠帝实录》建文四年条；《明史·兵志》及相关列传；边镇、地方志或对方材料。限度：「靖难之役：朱棣被拦在南京对面的长江边」借《惠帝实录》建文四年条；《明史·兵志》及相关列传；边镇、地方志或对方材料获得可检索的时间坐标；跨区域行政、资源或军政网络里朝廷与官员的路线、人数与后果若无其他材料相证，均保留为不可见的部分。
  \item[\eventanchor{ML-1402-009}{1402（建文四年）七月3日：靖难之役：布政司侍郎陈选投奔朱棣，起义军渡过长江}] 明与地方反抗、流动作战集团；朝廷与官员、军队与卫所、地方社会与精英、流动与反抗集团；地方或省域；未具城镇者不补造路线。发生在地方或省域的「靖难之役：布政司侍郎陈选投奔朱棣，起义军渡过长江」牵动朝廷与官员、军队与卫所、地方社会与精英、流动与反抗集团；已知材料让我们看见这一时点，却没有把准备行动、地方资源与决策意图完整连起来，因果只能保留为待核问题。 「靖难之役：布政司侍郎陈选投奔朱棣，起义军渡过长江」并未结束朝廷与官员、军队与卫所、地方社会与精英、流动与反抗集团在地方或省域的压力，而是把问题转到后续防务、运输或行政处置；只有后续记录能说明它实际造成何种改变。 材料：《惠帝实录》建文四年条；《明史·兵志》及相关列传；边镇、地方志或对方材料。限度：《惠帝实录》建文四年条；《明史·兵志》及相关列传；边镇、地方志或对方材料为「靖难之役：布政司侍郎陈选投奔朱棣，起义军渡过长江」留下编年入口，却未必保留地方或省域中朝廷与官员、军队与卫所、地方社会与精英、流动与反抗集团的完整视角；本条因此不据之裁定人数、责任或动机。
  \item[\eventanchor{ML-1402-010}{1402（建文四年）七月13日：靖难之战：朱晖打开南京金川门，不战而让朱棣入城；建文帝失踪，家人被监禁}] 明；朝廷与官员、军队与卫所；跨区域行政、资源或军政网络；不假定同步执行。发生在跨区域行政、资源或军政网络的「靖难之战：朱晖打开南京金川门，不战而让朱棣入城；建文帝失踪，家人被监禁」牵动朝廷与官员、军队与卫所；已知材料让我们看见这一时点，却没有把准备行动、地方资源与决策意图完整连起来，因果只能保留为待核问题。 「靖难之战：朱晖打开南京金川门，不战而让朱棣入城；建文帝失踪，家人被监禁」并未结束朝廷与官员、军队与卫所在跨区域行政、资源或军政网络的压力，而是把问题转到后续防务、运输或行政处置；只有后续记录能说明它实际造成何种改变。 材料：《惠帝实录》建文四年条；《明史·兵志》及相关列传；边镇、地方志或对方材料。限度：《惠帝实录》建文四年条；《明史·兵志》及相关列传；边镇、地方志或对方材料为「靖难之战：朱晖打开南京金川门，不战而让朱棣入城；建文帝失踪，家人被监禁」留下编年入口，却未必保留跨区域行政、资源或军政网络中朝廷与官员、军队与卫所的完整视角；本条因此不据之裁定人数、责任或动机。
  \item[\eventanchor{ML-1402-011}{1402（建文四年）七月17日：朱棣即位，是为永乐皇帝}] 明；朝廷与官员；跨区域行政、资源或军政网络；不假定同步执行。把“朱棣即位，是为永乐皇帝”放回1402（建文四年）七月17日，《惠帝实录》建文四年条；《明史·本纪》及相关列传；诏令、奏疏或会典版本只让我们看见朝廷与官员在跨区域的行政、资源与军政网络的一段处置链。 读“朱棣即位，是为永乐皇帝”的余波时，应转向1402（建文四年）七月17日之后的任命、批答或地方反应，不能把后出的道德评价倒灌进《惠帝实录》建文四年条；《明史·本纪》及相关列传；诏令、奏疏或会典版本所记的时点。 材料：《惠帝实录》建文四年条；《明史·本纪》及相关列传；诏令、奏疏或会典版本。限度：《惠帝实录》建文四年条；《明史·本纪》及相关列传；诏令、奏疏或会典版本使「朱棣即位，是为永乐皇帝」的时间位置可追查；至于朝廷与官员在跨区域行政、资源或军政网络的命令传递、路线和损失，必须保持待核而不能随编年补写。
  \item[\eventanchor{ML-1402-012}{1402（建文四年）九月：永乐皇帝委托编写《永乐百科全书》}] 明；朝廷与官员；跨区域行政、资源或军政网络；不假定同步执行。制度、教育或知识传播的需要推动了相关行动或文本形成。 它提供制度和传播史的锚点，不能据此推断社会整体接受。 材料：《惠帝实录》建文四年条；《明史·选举志》或《艺文志》；文集、碑刻或书目材料。限度：「永乐皇帝委托编写《永乐百科全书》」虽见于《惠帝实录》建文四年条；《明史·选举志》或《艺文志》；文集、碑刻或书目材料，但材料形成者未必观察到跨区域行政、资源或军政网络的全部过程；朝廷与官员的行动范围和受影响群体不能由此单独决定。
  \item[\eventanchor{EM-1403-YONGLE-DADIAN-COMMISSION}{1403（永乐元年）七月：朝廷命解缙等编纂大型类书，后称《永乐大典》}] 明；朝廷与官员；跨区域行政、资源或军政网络；不假定同步执行。新政权需吸纳士人、整理经史与展示文化正统；内阁文臣获得组织工程的机会 形成跨门类文献汇编项目；其抄写、贮藏和后世散佚另构成一条物质史 材料：《明太宗实录》永乐元年条；《明史·解缙传》；《永乐大典》现存卷册题记。限度：诏修及主编群体可靠；成书规模和保存本不能直接说明知识的社会传播范围
  \item[\eventanchor{ML-1403-001}{1403（永乐元年）二月：永乐皇帝封北平为“北都” 北京}] 明；朝廷与官员；跨区域行政、资源或军政网络；不假定同步执行。把“永乐皇帝封北平为“北都” 北京”放回1403（永乐元年）二月，《太宗实录》永乐元年条；《明史·本纪》及相关列传；诏令、奏疏或会典版本只让我们看见朝廷与官员在跨区域的行政、资源与军政网络的一段处置链。 读“永乐皇帝封北平为“北都” 北京”的余波时，应转向1403（永乐元年）二月之后的任命、批答或地方反应，不能把后出的道德评价倒灌进《太宗实录》永乐元年条；《明史·本纪》及相关列传；诏令、奏疏或会典版本所记的时点。 材料：《太宗实录》永乐元年条；《明史·本纪》及相关列传；诏令、奏疏或会典版本。限度：「永乐皇帝封北平为“北都” 北京」借《太宗实录》永乐元年条；《明史·本纪》及相关列传；诏令、奏疏或会典版本获得可检索的时间坐标；跨区域行政、资源或军政网络里朝廷与官员的路线、人数与后果若无其他材料相证，均保留为不可见的部分。
  \item[\eventanchor{ML-1403-002}{1403（永乐元年）四月：永乐皇帝在大宁附近安置了忠诚的乌里安凯}] 明与蒙古诸部；朝廷与官员、边地政权与社群；边地接触区；不把族名或部名当固定疆界。《太宗实录》永乐元年条；《明史·本纪》及相关列传；诏令、奏疏或会典版本的1403（永乐元年）四月条并未写成宫廷传记；它所确证的是朝廷与官员、边地政权与社群围绕“永乐皇帝在大宁附近安置了忠诚的乌里安凯”在账本所能辨认的空间尺度的行动痕迹。 “永乐皇帝在大宁附近安置了忠诚的乌里安凯”在1403（永乐元年）四月之后的人事、诏令或交涉应另随可见文书追索；《太宗实录》永乐元年条；《明史·本纪》及相关列传；诏令、奏疏或会典版本并不能独自裁定动机、密议和受影响者经验。 材料：《太宗实录》永乐元年条；《明史·本纪》及相关列传；诏令、奏疏或会典版本。限度：关于「永乐皇帝在大宁附近安置了忠诚的乌里安凯」，《太宗实录》永乐元年条；《明史·本纪》及相关列传；诏令、奏疏或会典版本所能确证的是条目所示的时点和行动；边地接触区内朝廷与官员、边地政权与社群的规模、意图及社会后果需要异源材料才能判断。
  \item[\eventanchor{ML-1403-003}{1403（永乐元年）九月4日：宝藏远航：下令建造200艘“远洋运输船”}] 明；朝廷与官员；跨区域行政、资源或军政网络；不假定同步执行。“宝藏远航：下令建造200艘“远洋运输船””在1403（永乐元年）九月4日进入《太宗实录》永乐元年条；《明史·本纪》及相关列传；诏令、奏疏或会典版本的叙述，因而可将朝廷与官员与跨区域的行政、资源与军政网络一并放回核验范围。 关于“宝藏远航：下令建造200艘“远洋运输船””的实际执行与代价，须从1403（永乐元年）九月4日之后的命令和回应另行检验；一份《太宗实录》永乐元年条；《明史·本纪》及相关列传；诏令、奏疏或会典版本的编年叙述并不足以代替它们。 材料：《太宗实录》永乐元年条；《明史·本纪》及相关列传；诏令、奏疏或会典版本。限度：《太宗实录》永乐元年条；《明史·本纪》及相关列传；诏令、奏疏或会典版本记录「宝藏远航：下令建造200艘“远洋运输船”」的方式限定了可说范围：可追踪1403（永乐元年）九月4日的行动节点，不能凭此还原跨区域行政、资源或军政网络中朝廷与官员的每一处部署。
  \item[\eventanchor{ML-1403-004}{1403（永乐元年）十二月：永乐皇帝设立建州卫}] 明；朝廷与官员、边地政权与社群；边地接触区；不把族名或部名当固定疆界。朝廷与官员、边地政权与社群在账本所能辨认的空间尺度的行动因“永乐皇帝设立建州卫”留有1403（永乐元年）十二月的《太宗实录》永乐元年条；《明史·本纪》及相关列传；诏令、奏疏或会典版本记录；这比概括性的权力动因更可检验。 “永乐皇帝设立建州卫”令后续核验获得1403（永乐元年）十二月这一时间坐标；《太宗实录》永乐元年条；《明史·本纪》及相关列传；诏令、奏疏或会典版本所见的行动者与空间不能推成完整权力因果或统一地方经验。 材料：《太宗实录》永乐元年条；《明史·本纪》及相关列传；诏令、奏疏或会典版本。限度：对「永乐皇帝设立建州卫」的《太宗实录》永乐元年条；《明史·本纪》及相关列传；诏令、奏疏或会典版本记录可支持年序核验；它不足以替代边地接触区的地方材料，故不把朝廷与官员、边地政权与社群的战报、奏报或叙事当成无争议结果。
  \item[\eventanchor{ML-1403-005}{1403（永乐元年）：日本对明朝中国的使团：足利义光派遣使者前往明朝，宣称自己是“您的臣民，日本国王”，并获得贸易特权}] 明与海上、朝贡相关政权；朝廷与官员、财政与商业行动者；账本可辨空间；地点细节仍须回到所列材料。朝贡名义、边境安全与港口贸易的利益使交涉持续发生。 它为后续册封、互市或海防安排留下文书与跨政权比较线索。 材料：《太宗实录》永乐元年条；《明史·外国传》；《朝鲜王朝实录》、琉球《历代宝案》或海外档案。限度：《太宗实录》永乐元年条；《明史·外国传》；《朝鲜王朝实录》、琉球《历代宝案》或海外档案使「日本对明朝中国的使团：足利义光派遣使者前往明朝，宣称自己是“您的臣民，日本国王”，并获得贸易特权」的时间位置可追查；至于朝廷与官员、财政与商业行动者在账本可辨空间的命令传递、路线和损失，必须保持待核而不能随编年补写。
  \item[\eventanchor{EM-1404-CROWN-PRINCE}{1404（永乐二年）：朱高炽被立为皇太子，继承竞争暂获制度性裁定}] 明；朝廷与官员；跨区域行政、资源或军政网络；不假定同步执行。朱高炽、朱高煦各有不同的官僚和军功支持；永乐须固定继承秩序 太子在南京监国和文官网络中积累经验；兄弟竞争仍在永乐晚年反复出现 材料：《明太宗实录》永乐二年条；《明史·仁宗本纪》与《汉王高煦传》；杨士奇《东里文集》东宫奏议。限度：册立可靠；太子与汉王的个人动机不可由后出党争叙事简单化
  \item[\eventanchor{ML-1404-001}{1404（永乐二年）三月：宝藏远航：下令建造50艘“海船”}] 明；朝廷与官员；跨区域行政、资源或军政网络；不假定同步执行。《太宗实录》永乐二年条；《明史·本纪》及相关列传；诏令、奏疏或会典版本的1404（永乐二年）三月条并未写成宫廷传记；它所确证的是朝廷与官员围绕“宝藏远航：下令建造50艘“海船””在跨区域的行政、资源与军政网络的行动痕迹。 “宝藏远航：下令建造50艘“海船””在1404（永乐二年）三月之后的人事、诏令或交涉应另随可见文书追索；《太宗实录》永乐二年条；《明史·本纪》及相关列传；诏令、奏疏或会典版本并不能独自裁定动机、密议和受影响者经验。 材料：《太宗实录》永乐二年条；《明史·本纪》及相关列传；诏令、奏疏或会典版本。限度：关于「宝藏远航：下令建造50艘“海船”」，《太宗实录》永乐二年条；《明史·本纪》及相关列传；诏令、奏疏或会典版本所能确证的是条目所示的时点和行动；跨区域行政、资源或军政网络内朝廷与官员的规模、意图及社会后果需要异源材料才能判断。
  \item[\eventanchor{ML-1404-002}{1404（永乐二年）七月：卡拉德尔的恩克帖木儿被明朝廷授予亲王称号}] 明；朝廷与官员；跨区域行政、资源或军政网络；不假定同步执行。“卡拉德尔的恩克帖木儿被明朝廷授予亲王称号”在1404（永乐二年）七月进入《太宗实录》永乐二年条；《明史·本纪》及相关列传；诏令、奏疏或会典版本的叙述，因而可将朝廷与官员与跨区域的行政、资源与军政网络一并放回核验范围。 关于“卡拉德尔的恩克帖木儿被明朝廷授予亲王称号”的实际执行与代价，须从1404（永乐二年）七月之后的命令和回应另行检验；一份《太宗实录》永乐二年条；《明史·本纪》及相关列传；诏令、奏疏或会典版本的编年叙述并不足以代替它们。 材料：《太宗实录》永乐二年条；《明史·本纪》及相关列传；诏令、奏疏或会典版本。限度：以《太宗实录》永乐二年条；《明史·本纪》及相关列传；诏令、奏疏或会典版本核对「卡拉德尔的恩克帖木儿被明朝廷授予亲王称号」时，只能确认其在1404（永乐二年）七月的叙述位置；朝廷与官员在跨区域行政、资源或军政网络的具体经验仍需由可定位的异类材料补证。
  \item[\eventanchor{ML-1404-003}{1404（永乐二年）十月：陈添平抵达南京，请求明朝恢复陈朝王位}] 明；朝廷与官员；跨区域行政、资源或军政网络；不假定同步执行。朝廷与官员在跨区域的行政、资源与军政网络的行动因“陈添平抵达南京，请求明朝恢复陈朝王位”留有1404（永乐二年）十月的《太宗实录》永乐二年条；《明史·本纪》及相关列传；诏令、奏疏或会典版本记录；这比概括性的权力动因更可检验。 “陈添平抵达南京，请求明朝恢复陈朝王位”令后续核验获得1404（永乐二年）十月这一时间坐标；《太宗实录》永乐二年条；《明史·本纪》及相关列传；诏令、奏疏或会典版本所见的行动者与空间不能推成完整权力因果或统一地方经验。 材料：《太宗实录》永乐二年条；《明史·本纪》及相关列传；诏令、奏疏或会典版本。限度：「陈添平抵达南京，请求明朝恢复陈朝王位」可由《太宗实录》永乐二年条；《明史·本纪》及相关列传；诏令、奏疏或会典版本定位到1404（永乐二年）十月，但这些叙述未完整保存跨区域行政、资源或军政网络的执行与受影响者经验，不能把沉默写成事实。
  \item[\eventanchor{ML-1404-004}{1404（永乐二年）十二月：帖木儿入侵明朝，但在途中去世}] 明；朝廷与官员；跨区域行政、资源或军政网络；不假定同步执行。《太宗实录》永乐二年条；《明史·本纪》及相关列传；诏令、奏疏或会典版本为1404（永乐二年）十二月的“帖木儿入侵明朝，但在途中去世”保留了一个窄而具体的观察面：朝廷与官员在跨区域的行政、资源与军政网络的行政动作。 “帖木儿入侵明朝，但在途中去世”留给后续政务的是一条从1404（永乐二年）十二月继续核对的文书次序；不能凭《太宗实录》永乐二年条；《明史·本纪》及相关列传；诏令、奏疏或会典版本判断谁意图何在或哪一派应负责任。 材料：《太宗实录》永乐二年条；《明史·本纪》及相关列传；诏令、奏疏或会典版本。限度：「帖木儿入侵明朝，但在途中去世」在1404（永乐二年）十二月的出现可回查《太宗实录》永乐二年条；《明史·本纪》及相关列传；诏令、奏疏或会典版本；它们不等于跨区域行政、资源或军政网络的全景记录，朝廷与官员未被记下的选择与损失不作推定。
  \item[\eventanchor{ML-1404-005}{1404（永乐二年）：恢复全国科举考试}] 明；朝廷与官员、宗教与文化行动者；跨区域行政、资源或军政网络；不假定同步执行。制度、教育或知识传播的需要推动了相关行动或文本形成。 它提供制度和传播史的锚点，不能据此推断社会整体接受。 材料：《太宗实录》永乐二年条；《明史·选举志》或《艺文志》；文集、碑刻或书目材料。限度：《太宗实录》永乐二年条；《明史·选举志》或《艺文志》；文集、碑刻或书目材料只能把「恢复全国科举考试」置入1404（永乐二年）的年序；对于跨区域行政、资源或军政网络中朝廷与官员、宗教与文化行动者的路线、人数和动机，仍须另据地方或对方材料核验。
  \item[\eventanchor{ML-1404-006}{1404（永乐二年）：山西万户搬迁北京}] 明；朝廷与官员；跨区域行政、资源或军政网络；不假定同步执行。“山西万户搬迁北京”的可证起点不是派系标签，而是1404（永乐二年）的《太宗实录》永乐二年条；《明史·本纪》及相关列传；诏令、奏疏或会典版本对朝廷与官员在跨区域的行政、资源与军政网络所作的登记。 “山西万户搬迁北京”在1404（永乐二年）之后的人事、诏令或交涉应另随可见文书追索；《太宗实录》永乐二年条；《明史·本纪》及相关列传；诏令、奏疏或会典版本并不能独自裁定动机、密议和受影响者经验。 材料：《太宗实录》永乐二年条；《明史·本纪》及相关列传；诏令、奏疏或会典版本。限度：《太宗实录》永乐二年条；《明史·本纪》及相关列传；诏令、奏疏或会典版本使「山西万户搬迁北京」的时间位置可追查；至于朝廷与官员在跨区域行政、资源或军政网络的命令传递、路线和损失，必须保持待核而不能随编年补写。
  \item[\eventanchor{EM-1405-ZHENGHE-FIRST-VOYAGE}{1405（永乐三年）六月：郑和率舰队出航，开始第一次西洋航行}] 明与印度洋港口政权；朝廷与官员；账本可辨空间；地点细节仍须回到所列材料。永乐朝借海上使团重建朝贡秩序并寻求海外情报和政治承认 官方舰队把东南港口、东南亚和印度洋使节网络连接起来；航行的财政和地方劳役成本需要另证 材料：《明太宗实录》永乐三年条；《明史·郑和传》；马欢《瀛涯胜览》相关航程。限度：出航可靠；船数、最大船型和贸易规模须区分官修记载、后出传说与物证
  \item[\eventanchor{ML-1405-001}{1405（永乐三年）七月11日：宝船航行：郑和率领27800人乘坐255艘船从南京出发，其中62艘为宝船，“载有圣旨给西洋各国，并根据身份等级，向其国王赠送金锦、花纹丝绸、彩纱”。}] 明；朝廷与官员；账本可辨空间；地点细节仍须回到所列材料。“宝船航行：郑和率领27800人乘坐255艘船从南京出发，其中62艘为宝船，“载有圣旨给西洋各国，并根据身份等级，向其国王赠送金锦、花纹丝绸、彩纱””在1405（永乐三年）七月11日进入《太宗实录》永乐三年条；《明史·本纪》及相关列传；诏令、奏疏或会典版本的叙述，因而可将朝廷与官员、海上与海外政权与账本所能辨认的空间尺度一并放回核验范围。 关于“宝船航行：郑和率领27800人乘坐255艘船从南京出发，其中62艘为宝船，“载有圣旨给西洋各国，并根据身份等级，向其国王赠送金锦、花纹丝绸、彩纱””的实际执行与代价，须从1405（永乐三年）七月11日之后的命令和回应另行检验；一份《太宗实录》永乐三年条；《明史·本纪》及相关列传；诏令、奏疏或会典版本的编年叙述并不足以代替它们。 材料：《太宗实录》永乐三年条；《明史·本纪》及相关列传；诏令、奏疏或会典版本。限度：就「宝船航行：郑和率领27800人乘坐255艘船从南京出发，其中62艘为宝船，“载有圣旨给西洋各国，并根据身份等级，向其国王赠送金锦、花纹丝绸、彩纱”」而论，《太宗实录》永乐三年条；《明史·本纪》及相关列传；诏令、奏疏或会典版本提供的是1405（永乐三年）七月11日的记录线索；朝廷与官员在账本可辨空间的实际行动细节不能由这一材料链独自补足。
  \item[\eventanchor{ML-1405-002}{1405（永乐三年）：北京新宫殿建筑开工}] 明；朝廷与官员、财政与商业行动者；跨区域行政、资源或军政网络；不假定同步执行。朝廷与官员、财政与商业行动者在跨区域的行政、资源与军政网络的行动因“北京新宫殿建筑开工”留有1405（永乐三年）的《太宗实录》永乐三年条；《明史·本纪》及相关列传；诏令、奏疏或会典版本记录；这比概括性的权力动因更可检验。 “北京新宫殿建筑开工”令后续核验获得1405（永乐三年）这一时间坐标；《太宗实录》永乐三年条；《明史·本纪》及相关列传；诏令、奏疏或会典版本所见的行动者与空间不能推成完整权力因果或统一地方经验。 材料：《太宗实录》永乐三年条；《明史·本纪》及相关列传；诏令、奏疏或会典版本。限度：对「北京新宫殿建筑开工」的《太宗实录》永乐三年条；《明史·本纪》及相关列传；诏令、奏疏或会典版本记录可支持年序核验；它不足以替代跨区域行政、资源或军政网络的地方材料，故不把朝廷与官员、财政与商业行动者的战报、奏报或叙事当成无争议结果。
  \item[\eventanchor{ML-1405-003}{1405（永乐三年）：北方营建、漕运与军需征解的本年文书显示资源调动压力，报数、起运和实际到达不得混同。}] 明；朝廷与官员、军队与卫所、财政与商业行动者；跨区域行政、资源或军政网络；不假定同步执行。就1405（永乐三年）而言，“北方营建、漕运与军需征解的本年文书显示资源调动压力，报数、起运和实际到达不得混同”让朝廷与官员、军队与卫所、财政与商业行动者在跨区域的行政、资源与军政网络的资源安排进入《太宗实录》永乐三年条；《明会典》相关条目；地方志、黄册/契约/河工材料的可见范围。 “北方营建、漕运与军需征解的本年文书显示资源调动压力，报数、起运和实际到达不得混同”只能把1405（永乐三年）的线索接到地方志、契约、河工或仓储材料；《太宗实录》永乐三年条；《明会典》相关条目；地方志、黄册/契约/河工材料所载的中央条文不能直接界定执行范围。 材料：《太宗实录》永乐三年条；《明会典》相关条目；地方志、黄册/契约/河工材料。限度：桥接条目只标示可回查的年度问题与材料链；实录有中央选择性，崇祯期尤其需要奏疏、地方、朝鲜及满文材料交叉。
  \item[\eventanchor{EM-1406-ANNAM-INVASION}{1406（永乐四年）：明以恢复陈氏为名出兵安南，进攻胡季犛政权}] 明与大虞政权；朝廷与官员、军队与卫所；账本可辨空间；地点细节仍须回到所列材料。胡氏取代陈朝引发流亡者求援；永乐朝将宗藩名义与南方扩张结合 战争迅速占领主要城市，却把明军置于持续的地方抵抗和补给压力之中 材料：《明太宗实录》永乐四年条；《明史·安南传》；越南《大越史记全书》相关纪年。限度：出兵和政治名义可靠；当地支持度与战场损失不能由明方檄文推断
  \item[\eventanchor{ML-1406-001}{1406（永乐四年）四月4日：陈添平和他的明朝护卫在进入琅山时遭到伏击并被杀}] 明；朝廷与官员；跨区域行政、资源或军政网络；不假定同步执行。朝廷与官员在跨区域的行政、资源与军政网络的行动因“陈添平和他的明朝护卫在进入琅山时遭到伏击并被杀”留有1406（永乐四年）四月4日的《太宗实录》永乐四年条；《明史·本纪》及相关列传；诏令、奏疏或会典版本记录；这比概括性的权力动因更可检验。 “陈添平和他的明朝护卫在进入琅山时遭到伏击并被杀”令后续核验获得1406（永乐四年）四月4日这一时间坐标；《太宗实录》永乐四年条；《明史·本纪》及相关列传；诏令、奏疏或会典版本所见的行动者与空间不能推成完整权力因果或统一地方经验。 材料：《太宗实录》永乐四年条；《明史·本纪》及相关列传；诏令、奏疏或会典版本。限度：「陈添平和他的明朝护卫在进入琅山时遭到伏击并被杀」可由《太宗实录》永乐四年条；《明史·本纪》及相关列传；诏令、奏疏或会典版本定位到1406（永乐四年）四月4日，但这些叙述未完整保存跨区域行政、资源或军政网络的执行与受影响者经验，不能把沉默写成事实。
  \item[\eventanchor{ML-1406-003}{1406（永乐四年）十二月13日：明胡战争：明军占领达邦和升龙}] 明；朝廷与官员、军队与卫所；跨区域行政、资源或军政网络；不假定同步执行。对于明的「明胡战争：明军占领达邦和升龙」，跨区域行政、资源或军政网络与朝廷与官员、军队与卫所限定了行动的可行范围；史料未逐项记录前置命令和供给，不能以后来结果补造一个整齐的起因。 「明胡战争：明军占领达邦和升龙」之后，朝廷与官员、军队与卫所在跨区域行政、资源或军政网络的行动空间有所改写；这种变化须由随后文书和地方材料追踪，不能把初报数字直接当作最终后果。 材料：《太宗实录》永乐四年条；《明史·兵志》及相关列传；边镇、地方志或对方材料。限度：关于「明胡战争：明军占领达邦和升龙」，《太宗实录》永乐四年条；《明史·兵志》及相关列传；边镇、地方志或对方材料所能确证的是条目所示的时点和行动；跨区域行政、资源或军政网络内朝廷与官员、军队与卫所的规模、意图及社会后果需要异源材料才能判断。
  \item[\eventanchor{ML-1406-004}{1406（永乐四年）：宝藏航行：宝藏船队访问马六甲和爪哇，然后沿马六甲海峡前往阿鲁、萨穆德拉巴赛苏丹国和兰布里，那里的人民被描述为“非常诚实和真诚”，从那里出发 3 天到达安达曼群岛，然后再用 8 天到达锡兰西海岸，那里的国王充满敌意。舰队启程前往被誉为“西洋大国”的卡利卡特}] 明与海上、朝贡相关政权；朝廷与官员；账本可辨空间；地点细节仍须回到所列材料。朝贡名义、边境安全与港口贸易的利益使交涉持续发生。 它为后续册封、互市或海防安排留下文书与跨政权比较线索。 材料：《太宗实录》永乐四年条；《明史·外国传》；《朝鲜王朝实录》、琉球《历代宝案》或海外档案。限度：《太宗实录》永乐四年条；《明史·外国传》；《朝鲜王朝实录》、琉球《历代宝案》或海外档案使「宝藏航行：宝藏船队访问马六甲和爪哇，然后沿马六甲海峡前往阿鲁、萨穆德拉巴赛苏丹国和兰布里，那里的人民被描述为“非常诚实和真诚”，从那里出发 3 天到达安达曼群岛，然后再用 8 天到达锡兰西海岸，那里的国王充满敌意。舰队启程前往被誉为“西洋大国”的卡利卡特」的时间位置可追查；至于朝廷与官员在账本可辨空间的命令传递、路线和损失，必须保持待核而不能随编年补写。
  \item[\eventanchor{EM-1407-JIAOZHI-PROVINCE}{1407（永乐五年）：明废大虞并设交趾承宣布政使司等机构}] 明与交趾地区；朝廷与官员；账本可辨空间；地点细节仍须回到所列材料。明军俘获胡氏父子后选择直接统治而非扶立陈氏后裔 官署、卫所和税役被移植到交趾；高成本占领激化地方抵抗并牵制南方军力 材料：《明太宗实录》永乐五年条；《明史·地理志》与《安南传》；越南《大越史记全书》。限度：设治诏令可靠；行政建置不能证明红河三角洲和山区已被有效控制
  \item[\eventanchor{EM-1407-ZHENGHE-SECOND-VOYAGE}{1407（永乐五年）：郑和第二次航行出发，护送各国使节返航}] 明与东南亚、南亚港口政权；朝廷与官员；跨区域行政、资源或军政网络；不假定同步执行。第一次航行带回多国使者，明廷以返送维持官方互惠与海上声望 航行由一次远征转为可重复的使节循环；港口补给和翻译者作用更加关键 材料：《明太宗实录》永乐五年条；《明史·郑和传》；费信《星槎胜览》。限度：航行存在可靠；访问港口的先后和个别使团身份需与碑刻、海外材料互校
  \item[\eventanchor{ML-1407-001}{1407（永乐五年）：宝藏航行：宝藏船队在巨港击败了陈祖义的海盗船队，并任命石进清为“统治当地原住民的大酋长”}] 明、后金（清）与辽东诸势力；朝廷与官员、边地政权与社群；账本可辨空间；地点细节仍须回到所列材料。「宝藏航行：宝藏船队在巨港击败了陈祖义的海盗船队，并任命石进清为“统治当地原住民的大酋长”」所涉朝廷与官员、边地政权与社群的行动受账本可辨空间的空间与资源约束；年序材料保存了节点，却没有保存足以裁定出兵、守备或支援动机的全过程。 「宝藏航行：宝藏船队在巨港击败了陈祖义的海盗船队，并任命石进清为“统治当地原住民的大酋长”」之后，账本可辨空间内朝廷与官员、边地政权与社群仍须处理由该行动留下的控制、通行或防务问题；本条不把战报中的胜负、斩获或伤亡扩大为无从核对的结论。 材料：《太宗实录》永乐五年条；《明史·兵志》及相关列传；边镇、地方志或对方材料。限度：《太宗实录》永乐五年条；《明史·兵志》及相关列传；边镇、地方志或对方材料使「宝藏航行：宝藏船队在巨港击败了陈祖义的海盗船队，并任命石进清为“统治当地原住民的大酋长”」的时间位置可追查；至于朝廷与官员、边地政权与社群在账本可辨空间的命令传递、路线和损失，必须保持待核而不能随编年补写。
  \item[\eventanchor{ML-1407-002}{1407（永乐五年）三月13日：明—大五（胡朝）战争：胡曲里对明军的反攻失败}] 明；朝廷与官员、军队与卫所；跨区域行政、资源或军政网络；不假定同步执行。「明—大五（胡朝）战争：胡曲里对明军的反攻失败」把朝廷与官员、军队与卫所带到跨区域行政、资源或军政网络的具体关口；现存年条只留下这一行动节点，前期兵力、粮道和地方协力的次序未逐项保存，不能倒推为单一动因。 「明—大五（胡朝）战争：胡曲里对明军的反攻失败」改变了朝廷与官员、军队与卫所处理跨区域行政、资源或军政网络事务的顺序与代价；后续影响须由连贯记录显示，不能从孤立军报推出可量化的整体结局。 材料：《太宗实录》永乐五年条；《明史·兵志》及相关列传；边镇、地方志或对方材料。限度：「明—大五（胡朝）战争：胡曲里对明军的反攻失败」虽见于《太宗实录》永乐五年条；《明史·兵志》及相关列传；边镇、地方志或对方材料，但材料形成者未必观察到跨区域行政、资源或军政网络的全部过程；朝廷与官员、军队与卫所的行动范围和受影响群体不能由此单独决定。
  \item[\eventanchor{ML-1407-003}{1407（永乐五年）四月：五世噶玛巴喇嘛德新谢巴抵达南京举行宗教仪式}] 明；朝廷与官员、宗教与文化行动者；跨区域行政、资源或军政网络；不假定同步执行。“五世噶玛巴喇嘛德新谢巴抵达南京举行宗教仪式”的可证起点不是派系标签，而是1407（永乐五年）四月的《太宗实录》永乐五年条；《明史·本纪》及相关列传；诏令、奏疏或会典版本对朝廷与官员、宗教与文化行动者在跨区域的行政、资源与军政网络所作的登记。 “五世噶玛巴喇嘛德新谢巴抵达南京举行宗教仪式”在1407（永乐五年）四月之后的人事、诏令或交涉应另随可见文书追索；《太宗实录》永乐五年条；《明史·本纪》及相关列传；诏令、奏疏或会典版本并不能独自裁定动机、密议和受影响者经验。 材料：《太宗实录》永乐五年条；《明史·本纪》及相关列传；诏令、奏疏或会典版本。限度：就「五世噶玛巴喇嘛德新谢巴抵达南京举行宗教仪式」而论，《太宗实录》永乐五年条；《明史·本纪》及相关列传；诏令、奏疏或会典版本提供的是1407（永乐五年）四月的记录线索；朝廷与官员、宗教与文化行动者在跨区域行政、资源或军政网络的实际行动细节不能由这一材料链独自补足。
  \item[\eventanchor{ML-1407-004}{1407（永乐五年）六月16日：明—大五（胡朝）战争：胡贵利父子被俘，送往南京}] 明；朝廷与官员、军队与卫所；跨区域行政、资源或军政网络；不假定同步执行。从跨区域行政、资源或军政网络看「明—大五（胡朝）战争：胡贵利父子被俘，送往南京」，朝廷与官员、军队与卫所的行动连接着当年的军政网络；资料没有充分说明何种先行安排直接促成它，不能把编年相邻当作因果证明。 「明—大五（胡朝）战争：胡贵利父子被俘，送往南京」并未结束朝廷与官员、军队与卫所在跨区域行政、资源或军政网络的压力，而是把问题转到后续防务、运输或行政处置；只有后续记录能说明它实际造成何种改变。 材料：《太宗实录》永乐五年条；《明史·兵志》及相关列传；边镇、地方志或对方材料。限度：《太宗实录》永乐五年条；《明史·兵志》及相关列传；边镇、地方志或对方材料为「明—大五（胡朝）战争：胡贵利父子被俘，送往南京」留下编年入口，却未必保留跨区域行政、资源或军政网络中朝廷与官员、军队与卫所的完整视角；本条因此不据之裁定人数、责任或动机。
  \item[\eventanchor{ML-1407-005}{1407（永乐五年）七月5日：中国第四次统治越南：永乐皇帝宣布交趾正式并入明朝}] 明；朝廷与官员；账本可辨空间；地点细节仍须回到所列材料。朝贡名义、边境安全与港口贸易的利益使交涉持续发生。 它为后续册封、互市或海防安排留下文书与跨政权比较线索。 材料：《太宗实录》永乐五年条；《明史·外国传》；《朝鲜王朝实录》、琉球《历代宝案》或海外档案。限度：「中国第四次统治越南：永乐皇帝宣布交趾正式并入明朝」借《太宗实录》永乐五年条；《明史·外国传》；《朝鲜王朝实录》、琉球《历代宝案》或海外档案获得可检索的时间坐标；账本可辨空间里朝廷与官员的路线、人数与后果若无其他材料相证，均保留为不可见的部分。
  \item[\eventanchor{ML-1407-006}{1407（永乐五年）十月2日：宝藏之旅：中国宝藏船队返回南京}] 明；朝廷与官员；跨区域行政、资源或军政网络；不假定同步执行。《太宗实录》永乐五年条；《明史·本纪》及相关列传；诏令、奏疏或会典版本把1407（永乐五年）十月2日的“宝藏之旅：中国宝藏船队返回南京”编入年序；能直接追索的是朝廷与官员在跨区域的行政、资源与军政网络所面对的那一项处置。 关于“宝藏之旅：中国宝藏船队返回南京”的实际执行与代价，须从1407（永乐五年）十月2日之后的命令和回应另行检验；一份《太宗实录》永乐五年条；《明史·本纪》及相关列传；诏令、奏疏或会典版本的编年叙述并不足以代替它们。 材料：《太宗实录》永乐五年条；《明史·本纪》及相关列传；诏令、奏疏或会典版本。限度：关于「宝藏之旅：中国宝藏船队返回南京」，《太宗实录》永乐五年条；《明史·本纪》及相关列传；诏令、奏疏或会典版本所能确证的是条目所示的时点和行动；跨区域行政、资源或军政网络内朝廷与官员的规模、意图及社会后果需要异源材料才能判断。
  \item[\eventanchor{ML-1407-007}{1407（永乐五年）十月5日：宝藏远航：王浩奉命改装249艘“海上运输船”，为“驻西洋各国使馆做准备”}] 明；朝廷与官员；账本可辨空间；地点细节仍须回到所列材料。沿着1407（永乐五年）十月5日的《太宗实录》永乐五年条；《明史·本纪》及相关列传；诏令、奏疏或会典版本回看，“宝藏远航：王浩奉命改装249艘“海上运输船”，为“驻西洋各国使馆做准备””首先是朝廷与官员、海上与海外政权在账本所能辨认的空间尺度留下的记录，而非后来人物评判的注脚。 “宝藏远航：王浩奉命改装249艘“海上运输船”，为“驻西洋各国使馆做准备””令后续核验获得1407（永乐五年）十月5日这一时间坐标；《太宗实录》永乐五年条；《明史·本纪》及相关列传；诏令、奏疏或会典版本所见的行动者与空间不能推成完整权力因果或统一地方经验。 材料：《太宗实录》永乐五年条；《明史·本纪》及相关列传；诏令、奏疏或会典版本。限度：《太宗实录》永乐五年条；《明史·本纪》及相关列传；诏令、奏疏或会典版本记录「宝藏远航：王浩奉命改装249艘“海上运输船”，为“驻西洋各国使馆做准备”」的方式限定了可说范围：可追踪1407（永乐五年）十月5日的行动节点，不能凭此还原账本可辨空间中朝廷与官员的每一处部署。
  \item[\eventanchor{ML-1407-008}{1407（永乐五年）十月30日：宝藏之旅：一位太监总督带着一封给占婆国王的御书启程}] 明与海上、朝贡相关政权；朝廷与官员；账本可辨空间；地点细节仍须回到所列材料。朝贡名义、边境安全与港口贸易的利益使交涉持续发生。 它为后续册封、互市或海防安排留下文书与跨政权比较线索。 材料：《太宗实录》永乐五年条；《明史·外国传》；《朝鲜王朝实录》、琉球《历代宝案》或海外档案。限度：《太宗实录》永乐五年条；《明史·外国传》；《朝鲜王朝实录》、琉球《历代宝案》或海外档案为「宝藏之旅：一位太监总督带着一封给占婆国王的御书启程」留下编年入口，却未必保留账本可辨空间中朝廷与官员的完整视角；本条因此不据之裁定人数、责任或动机。
  \item[\eventanchor{ML-1407-009}{1407（永乐五年）：郑和下西洋：郑和船队由249艘船组成，航线与第一次下西洋相似，增加了嘉义勒、阿博巴丹、甘巴里、奎隆、科钦等地}] 明；朝廷与官员；账本可辨空间；地点细节仍须回到所列材料。1407（永乐五年）的材料链以《太宗实录》永乐五年条；《明史·本纪》及相关列传；诏令、奏疏或会典版本为轴，所能落实的是“郑和下西洋：郑和船队由249艘船组成，航线与第一次下西洋相似，增加了嘉义勒、阿博巴丹、甘巴里、奎隆、科钦等地”与朝廷与官员、海上与海外政权在账本所能辨认的空间尺度的相遇。 读“郑和下西洋：郑和船队由249艘船组成，航线与第一次下西洋相似，增加了嘉义勒、阿博巴丹、甘巴里、奎隆、科钦等地”的余波时，应转向1407（永乐五年）之后的任命、批答或地方反应，不能把后出的道德评价倒灌进《太宗实录》永乐五年条；《明史·本纪》及相关列传；诏令、奏疏或会典版本所记的时点。 材料：《太宗实录》永乐五年条；《明史·本纪》及相关列传；诏令、奏疏或会典版本。限度：「郑和下西洋：郑和船队由249艘船组成，航线与第一次下西洋相似，增加了嘉义勒、阿博巴丹、甘巴里、奎隆、科钦等地」借《太宗实录》永乐五年条；《明史·本纪》及相关列传；诏令、奏疏或会典版本获得可检索的时间坐标；账本可辨空间里朝廷与官员的路线、人数与后果若无其他材料相证，均保留为不可见的部分。
  \item[\eventanchor{ML-1407-010}{1407（永乐五年）十二月：《永乐百科全书》完成}] 明；朝廷与官员；跨区域行政、资源或军政网络；不假定同步执行。制度、教育或知识传播的需要推动了相关行动或文本形成。 它提供制度和传播史的锚点，不能据此推断社会整体接受。 材料：《太宗实录》永乐五年条；《明史·选举志》或《艺文志》；文集、碑刻或书目材料。限度：以《太宗实录》永乐五年条；《明史·选举志》或《艺文志》；文集、碑刻或书目材料核对「《永乐百科全书》完成」时，只能确认其在1407（永乐五年）十二月的叙述位置；朝廷与官员在跨区域行政、资源或军政网络的具体经验仍需由可定位的异类材料补证。
  \item[\eventanchor{ML-1407-011}{1407（永乐五年）：明朝大炮中添加了铁木填料，提高了其效能。}] 明；朝廷与官员；跨区域行政、资源或军政网络；不假定同步执行。制度、教育或知识传播的需要推动了相关行动或文本形成。 它提供制度和传播史的锚点，不能据此推断社会整体接受。 材料：《太宗实录》永乐五年条；《明史·选举志》或《艺文志》；文集、碑刻或书目材料。限度：对「明朝大炮中添加了铁木填料，提高了其效能」的《太宗实录》永乐五年条；《明史·选举志》或《艺文志》；文集、碑刻或书目材料记录可支持年序核验；它不足以替代跨区域行政、资源或军政网络的地方材料，故不把朝廷与官员的战报、奏报或叙事当成无争议结果。
  \item[\eventanchor{EM-1408-YONGLE-DADIAN-COMPLETED}{1408（永乐六年）：《永乐大典》进呈，类书编纂工程完成一阶段}] 明；朝廷与官员、财政与商业行动者；跨区域行政、资源或军政网络；不假定同步执行。朝廷投入大批儒臣与抄手完成分类汇编，以文化工程巩固永乐正统 文献汇编成为内府收藏与后世辑佚的重要媒介；政治整肃中的编者命运也显示工程依赖宫廷权力 材料：《明太宗实录》永乐六年条；《明史·解缙传》；现存《永乐大典》卷端题记。限度：进呈可靠；“收尽天下书”是工程宣称，不能以之推定文献无遗漏
  \item[\eventanchor{ML-1408-001}{1408（永乐六年）二月14日：宝藏远航：南京工部下令建造48艘宝船}] 明；朝廷与官员、财政与商业行动者；跨区域行政、资源或军政网络；不假定同步执行。就1408（永乐六年）二月14日而言，《太宗实录》永乐六年条；《明史·本纪》及相关列传；诏令、奏疏或会典版本记录“宝藏远航：南京工部下令建造48艘宝船”时同时划出朝廷与官员、财政与商业行动者与跨区域的行政、资源与军政网络这两个不能省略的条件。 读“宝藏远航：南京工部下令建造48艘宝船”的余波时，应转向1408（永乐六年）二月14日之后的任命、批答或地方反应，不能把后出的道德评价倒灌进《太宗实录》永乐六年条；《明史·本纪》及相关列传；诏令、奏疏或会典版本所记的时点。 材料：《太宗实录》永乐六年条；《明史·本纪》及相关列传；诏令、奏疏或会典版本。限度：「宝藏远航：南京工部下令建造48艘宝船」虽见于《太宗实录》永乐六年条；《明史·本纪》及相关列传；诏令、奏疏或会典版本，但材料形成者未必观察到跨区域行政、资源或军政网络的全部过程；朝廷与官员、财政与商业行动者的行动范围和受影响群体不能由此单独决定。
  \item[\eventanchor{ML-1408-002}{1408（永乐六年）七月5日：中国第四次统治越南：明军夺取1360万吨大米；越南有235,900头牛、牲畜和大量物资}] 明；朝廷与官员、军队与卫所；跨区域行政、资源或军政网络；不假定同步执行。朝贡名义、边境安全与港口贸易的利益使交涉持续发生。 它为后续册封、互市或海防安排留下文书与跨政权比较线索。 材料：《太宗实录》永乐六年条；《明史·外国传》；《朝鲜王朝实录》、琉球《历代宝案》或海外档案。限度：「中国第四次统治越南：明军夺取1360万吨大米；越南有235,900头牛、牲畜和大量物资」借《太宗实录》永乐六年条；《明史·外国传》；《朝鲜王朝实录》、琉球《历代宝案》或海外档案获得可检索的时间坐标；跨区域行政、资源或军政网络里朝廷与官员、军队与卫所的路线、人数与后果若无其他材料相证，均保留为不可见的部分。
  \item[\eventanchor{ML-1408-003}{1408（永乐六年）九月：中国第四次统治越南：交趾的陈毅叛军}] 明与地方反抗、流动作战集团；朝廷与官员、军队与卫所、地方社会与精英、流动与反抗集团；账本可辨空间；地点细节仍须回到所列材料。朝贡名义、边境安全与港口贸易的利益使交涉持续发生。 它为后续册封、互市或海防安排留下文书与跨政权比较线索。 材料：《太宗实录》永乐六年条；《明史·外国传》；《朝鲜王朝实录》、琉球《历代宝案》或海外档案。限度：《太宗实录》永乐六年条；《明史·外国传》；《朝鲜王朝实录》、琉球《历代宝案》或海外档案为「中国第四次统治越南：交趾的陈毅叛军」留下编年入口，却未必保留账本可辨空间中朝廷与官员、军队与卫所、地方社会与精英、流动与反抗集团的完整视角；本条因此不据之裁定人数、责任或动机。
  \item[\eventanchor{ML-1409-001}{1409（永乐七年）一月：宝藏航次：第三次航次下达指令}] 明；朝廷与官员；跨区域行政、资源或军政网络；不假定同步执行。“宝藏航次：第三次航次下达指令”的可证起点不是派系标签，而是1409（永乐七年）一月的《太宗实录》永乐七年条；《明史·本纪》及相关列传；诏令、奏疏或会典版本对朝廷与官员在跨区域的行政、资源与军政网络所作的登记。 “宝藏航次：第三次航次下达指令”在1409（永乐七年）一月之后的人事、诏令或交涉应另随可见文书追索；《太宗实录》永乐七年条；《明史·本纪》及相关列传；诏令、奏疏或会典版本并不能独自裁定动机、密议和受影响者经验。 材料：《太宗实录》永乐七年条；《明史·本纪》及相关列传；诏令、奏疏或会典版本。限度：「宝藏航次：第三次航次下达指令」在1409（永乐七年）一月的出现可回查《太宗实录》永乐七年条；《明史·本纪》及相关列传；诏令、奏疏或会典版本；它们不等于跨区域行政、资源或军政网络的全景记录，朝廷与官员未被记下的选择与损失不作推定。
  \item[\eventanchor{ML-1409-002}{1409（永乐七年）二月15日：宝藏之旅：加勒三语铭文诞生}] 明；朝廷与官员；跨区域行政、资源或军政网络；不假定同步执行。在1409（永乐七年）二月15日，《太宗实录》永乐七年条；《明史·本纪》及相关列传；诏令、奏疏或会典版本记下“宝藏之旅：加勒三语铭文诞生”；这使朝廷与官员在跨区域的行政、资源与军政网络的处置留下可核的文书入口。 读“宝藏之旅：加勒三语铭文诞生”的余波时，应转向1409（永乐七年）二月15日之后的任命、批答或地方反应，不能把后出的道德评价倒灌进《太宗实录》永乐七年条；《明史·本纪》及相关列传；诏令、奏疏或会典版本所记的时点。 材料：《太宗实录》永乐七年条；《明史·本纪》及相关列传；诏令、奏疏或会典版本。限度：「宝藏之旅：加勒三语铭文诞生」可由《太宗实录》永乐七年条；《明史·本纪》及相关列传；诏令、奏疏或会典版本定位到1409（永乐七年）二月15日，但这些叙述未完整保存跨区域行政、资源或军政网络的执行与受影响者经验，不能把沉默写成事实。
  \item[\eventanchor{ML-1409-003}{1409（永乐七年）六月：卫拉特从明朝廷获得王子头衔}] 明与蒙古诸部；朝廷与官员、边地政权与社群；边地接触区；不把族名或部名当固定疆界。“卫拉特从明朝廷获得王子头衔”见于《太宗实录》永乐七年条；《明史·本纪》及相关列传；诏令、奏疏或会典版本的1409（永乐七年）六月条。它把朝廷与官员、边地政权与社群置于账本所能辨认的空间尺度，而不是抽象的宫廷舞台。 “卫拉特从明朝廷获得王子头衔”在1409（永乐七年）六月之后的人事、诏令或交涉应另随可见文书追索；《太宗实录》永乐七年条；《明史·本纪》及相关列传；诏令、奏疏或会典版本并不能独自裁定动机、密议和受影响者经验。 材料：《太宗实录》永乐七年条；《明史·本纪》及相关列传；诏令、奏疏或会典版本。限度：以《太宗实录》永乐七年条；《明史·本纪》及相关列传；诏令、奏疏或会典版本核对「卫拉特从明朝廷获得王子头衔」时，只能确认其在1409（永乐七年）六月的叙述位置；朝廷与官员、边地政权与社群在边地接触区的具体经验仍需由可定位的异类材料补证。
  \item[\eventanchor{ML-1409-004}{1409（永乐七年）夏：宝藏之旅：宝船队返回中国}] 明；朝廷与官员；跨区域行政、资源或军政网络；不假定同步执行。《太宗实录》永乐七年条；《明史·本纪》及相关列传；诏令、奏疏或会典版本把1409（永乐七年）夏的“宝藏之旅：宝船队返回中国”编入年序；能直接追索的是朝廷与官员在跨区域的行政、资源与军政网络所面对的那一项处置。 关于“宝藏之旅：宝船队返回中国”的实际执行与代价，须从1409（永乐七年）夏之后的命令和回应另行检验；一份《太宗实录》永乐七年条；《明史·本纪》及相关列传；诏令、奏疏或会典版本的编年叙述并不足以代替它们。 材料：《太宗实录》永乐七年条；《明史·本纪》及相关列传；诏令、奏疏或会典版本。限度：关于「宝藏之旅：宝船队返回中国」，《太宗实录》永乐七年条；《明史·本纪》及相关列传；诏令、奏疏或会典版本所能确证的是条目所示的时点和行动；跨区域行政、资源或军政网络内朝廷与官员的规模、意图及社会后果需要异源材料才能判断。
  \item[\eventanchor{ML-1409-005}{1409（永乐七年）九月23日：克鲁伦之战：明军被乌列帖木儿汗击败}] 明；朝廷与官员、军队与卫所；跨区域行政、资源或军政网络；不假定同步执行。围绕「克鲁伦之战：明军被乌列帖木儿汗击败」，明在跨区域行政、资源或军政网络面对的是朝廷与官员、军队与卫所的现场处置；材料确认行动进入编年，却不足以把此前调兵、补给或地方压力串成确定因果链。 发生「克鲁伦之战：明军被乌列帖木儿汗击败」后，跨区域行政、资源或军政网络仍是朝廷与官员、军队与卫所必须面对的行动场域；可见结果是问题被带入下一段年序，未见的伤亡、掠夺与居民去留不在本条擅断。 材料：《太宗实录》永乐七年条；《明史·兵志》及相关列传；边镇、地方志或对方材料。限度：《太宗实录》永乐七年条；《明史·兵志》及相关列传；边镇、地方志或对方材料记录「克鲁伦之战：明军被乌列帖木儿汗击败」的方式限定了可说范围：可追踪1409（永乐七年）九月23日的行动节点，不能凭此还原跨区域行政、资源或军政网络中朝廷与官员、军队与卫所的每一处部署。
  \item[\eventanchor{ML-1409-006}{1409（永乐七年）十月：郑和下西洋：率领27000人，按常规路线出发}] 明；朝廷与官员；账本可辨空间；地点细节仍须回到所列材料。对“郑和下西洋：率领27000人，按常规路线出发”的最小把握来自1409（永乐七年）十月的《太宗实录》永乐七年条；《明史·本纪》及相关列传；诏令、奏疏或会典版本：朝廷与官员、海上与海外政权在账本所能辨认的空间尺度被纳入同一条行政记录。 “郑和下西洋：率领27000人，按常规路线出发”留给后续政务的是一条从1409（永乐七年）十月继续核对的文书次序；不能凭《太宗实录》永乐七年条；《明史·本纪》及相关列传；诏令、奏疏或会典版本判断谁意图何在或哪一派应负责任。 材料：《太宗实录》永乐七年条；《明史·本纪》及相关列传；诏令、奏疏或会典版本。限度：《太宗实录》永乐七年条；《明史·本纪》及相关列传；诏令、奏疏或会典版本为「郑和下西洋：率领27000人，按常规路线出发」留下编年入口，却未必保留账本可辨空间中朝廷与官员的完整视角；本条因此不据之裁定人数、责任或动机。
  \item[\eventanchor{ML-1409-007}{1409（永乐七年）十二月：中国第四次统治越南：明朝军队攻占陈鹗，但陈贵康成为叛军领袖}] 明与地方反抗、流动作战集团；朝廷与官员、军队与卫所、地方社会与精英、流动与反抗集团；地方或省域；未具城镇者不补造路线。朝贡名义、边境安全与港口贸易的利益使交涉持续发生。 它为后续册封、互市或海防安排留下文书与跨政权比较线索。 材料：《太宗实录》永乐七年条；《明史·外国传》；《朝鲜王朝实录》、琉球《历代宝案》或海外档案。限度：「中国第四次统治越南：明朝军队攻占陈鹗，但陈贵康成为叛军领袖」借《太宗实录》永乐七年条；《明史·外国传》；《朝鲜王朝实录》、琉球《历代宝案》或海外档案获得可检索的时间坐标；地方或省域里朝廷与官员、军队与卫所、地方社会与精英、流动与反抗集团的路线、人数与后果若无其他材料相证，均保留为不可见的部分。
  \item[\eventanchor{EM-1410-FIRST-MONGOL-CAMPAIGN}{1410（永乐八年）：永乐帝亲征蒙古，北上击败本雅失里及相关部众}] 明与鞑靼、瓦剌诸部；朝廷与官员、军队与卫所、边地政权与社群；边地接触区；不把族名或部名当固定疆界。本雅失里处置明使及草原联盟变化挑战明廷北方威望 永乐以皇帝亲征重塑威慑；草原政治未被消除，明军远征补给负担持续扩大 材料：《明太宗实录》永乐八年条；《明史·成祖本纪》与《鞑靼传》；17世纪《蒙古源流》鞑靼、瓦剌叙事。限度：亲征和作战可靠；战果、敌军数字和草原部众归属须以多方材料审读
  \item[\eventanchor{ML-1410-001}{1410（永乐八年）六月15日：第一次蒙古战役：永乐皇帝在斡难河畔击败额列帖木儿汗}] 明与蒙古诸部；朝廷与官员、军队与卫所、边地政权与社群；边地接触区；不把族名或部名当固定疆界。在1410（永乐八年）六月15日，朝廷与官员、军队与卫所、边地政权与社群因“第一次蒙古战役：永乐皇帝在斡难河畔击败额列帖木儿汗”进入账本所能辨认的空间尺度的行政记录，材料链为《太宗实录》永乐八年条；《明会典》相关条目；地方志、黄册/契约/河工材料。 “第一次蒙古战役：永乐皇帝在斡难河畔击败额列帖木儿汗”在1410（永乐八年）六月15日形成的是资源调度的核验线，而非完成的财政结论；到达、拖欠与成本仍要借异类材料追踪。 材料：《太宗实录》永乐八年条；《明会典》相关条目；地方志、黄册/契约/河工材料。限度：就「第一次蒙古战役：永乐皇帝在斡难河畔击败额列帖木儿汗」而论，《太宗实录》永乐八年条；《明会典》相关条目；地方志、黄册/契约/河工材料提供的是1410（永乐八年）六月15日的记录线索；朝廷与官员、军队与卫所、边地政权与社群在边地接触区的实际行动细节不能由这一材料链独自补足。
  \item[\eventanchor{ML-1410-002}{1410（永乐八年）七月：第一次蒙古战役：明军击败大兴安岭以东的阿鲁台并撤至南京}] 明与蒙古诸部；朝廷与官员、军队与卫所、边地政权与社群；边地接触区；不把族名或部名当固定疆界。在边地接触区发生的「第一次蒙古战役：明军击败大兴安岭以东的阿鲁台并撤至南京」使朝廷与官员、军队与卫所、边地政权与社群必须作出即时选择；可见记录没有完整交代先前部署，故只把军政背景列为条件，不写成已经证实的原因。 「第一次蒙古战役：明军击败大兴安岭以东的阿鲁台并撤至南京」并未结束朝廷与官员、军队与卫所、边地政权与社群在边地接触区的压力，而是把问题转到后续防务、运输或行政处置；只有后续记录能说明它实际造成何种改变。 材料：《太宗实录》永乐八年条；《明史·兵志》及相关列传；边镇、地方志或对方材料。限度：《太宗实录》永乐八年条；《明史·兵志》及相关列传；边镇、地方志或对方材料为「第一次蒙古战役：明军击败大兴安岭以东的阿鲁台并撤至南京」留下编年入口，却未必保留边地接触区中朝廷与官员、军队与卫所、边地政权与社群的完整视角；本条因此不据之裁定人数、责任或动机。
  \item[\eventanchor{ML-1410-003}{1410（永乐八年）：明科特战争：宝藏船队登陆锡兰加勒并俘获甘波拉王国国王维杰亚巴胡六世}] 明；朝廷与官员、军队与卫所；跨区域行政、资源或军政网络；不假定同步执行。「明科特战争：宝藏船队登陆锡兰加勒并俘获甘波拉王国国王维杰亚巴胡六世」所涉朝廷与官员、军队与卫所的行动受跨区域行政、资源或军政网络的空间与资源约束；年序材料保存了节点，却没有保存足以裁定出兵、守备或支援动机的全过程。 「明科特战争：宝藏船队登陆锡兰加勒并俘获甘波拉王国国王维杰亚巴胡六世」之后，跨区域行政、资源或军政网络内朝廷与官员、军队与卫所仍须处理由该行动留下的控制、通行或防务问题；本条不把战报中的胜负、斩获或伤亡扩大为无从核对的结论。 材料：《太宗实录》永乐八年条；《明史·兵志》及相关列传；边镇、地方志或对方材料。限度：《太宗实录》永乐八年条；《明史·兵志》及相关列传；边镇、地方志或对方材料使「明科特战争：宝藏船队登陆锡兰加勒并俘获甘波拉王国国王维杰亚巴胡六世」的时间位置可追查；至于朝廷与官员、军队与卫所在跨区域行政、资源或军政网络的命令传递、路线和损失，必须保持待核而不能随编年补写。
  \item[\eventanchor{EM-1411-GRAND-CANAL}{1411（永乐九年）：朝廷大规模疏浚会通河并调整运河漕运工程}] 明；朝廷与官员、财政与商业行动者；跨区域行政、资源或军政网络；不假定同步执行。北京经营需要稳定的江南粮运；黄河、会通河淤塞限制旧有水路 漕粮北运条件得到改善；河工征发、治黄与运河维护成为长期财政和地方社会负担 材料：《明太宗实录》永乐九年河工条；《明史·河渠志》；会通河遗址调查与运河沿线碑刻。限度：工程命令可靠；通航时间、役夫数量和受益地区需分段核对
  \item[\eventanchor{ML-1411-001}{1411（永乐九年）七月：大运河疏浚重建工程启动}] 明；朝廷与官员、财政与商业行动者；跨区域行政、资源或军政网络；不假定同步执行。《太宗实录》永乐九年条；《明会典》相关条目；地方志、黄册/契约/河工材料所见的1411（永乐九年）七月事务是“大运河疏浚重建工程启动”；它使朝廷与官员、财政与商业行动者在跨区域的行政、资源与军政网络的负担与选择成为可核对象。 “大运河疏浚重建工程启动”只能把1411（永乐九年）七月的线索接到地方志、契约、河工或仓储材料；《太宗实录》永乐九年条；《明会典》相关条目；地方志、黄册/契约/河工材料所载的中央条文不能直接界定执行范围。 材料：《太宗实录》永乐九年条；《明会典》相关条目；地方志、黄册/契约/河工材料。限度：《太宗实录》永乐九年条；《明会典》相关条目；地方志、黄册/契约/河工材料只能把「大运河疏浚重建工程启动」置入1411（永乐九年）七月的年序；对于跨区域行政、资源或军政网络中朝廷与官员、财政与商业行动者的路线、人数和动机，仍须另据地方或对方材料核验。
  \item[\eventanchor{ML-1411-002}{1411（永乐九年）七月6日：宝藏之旅：宝船队返回南京}] 明；朝廷与官员；跨区域行政、资源或军政网络；不假定同步执行。把“宝藏之旅：宝船队返回南京”放回1411（永乐九年）七月6日，《太宗实录》永乐九年条；《明史·本纪》及相关列传；诏令、奏疏或会典版本只让我们看见朝廷与官员在跨区域的行政、资源与军政网络的一段处置链。 读“宝藏之旅：宝船队返回南京”的余波时，应转向1411（永乐九年）七月6日之后的任命、批答或地方反应，不能把后出的道德评价倒灌进《太宗实录》永乐九年条；《明史·本纪》及相关列传；诏令、奏疏或会典版本所记的时点。 材料：《太宗实录》永乐九年条；《明史·本纪》及相关列传；诏令、奏疏或会典版本。限度：以《太宗实录》永乐九年条；《明史·本纪》及相关列传；诏令、奏疏或会典版本核对「宝藏之旅：宝船队返回南京」时，只能确认其在1411（永乐九年）七月6日的叙述位置；朝廷与官员在跨区域行政、资源或军政网络的具体经验仍需由可定位的异类材料补证。
  \item[\eventanchor{ML-1411-003}{1411（永乐九年）：永乐皇帝派遣一失哈探索北满洲}] 明、后金（清）与辽东诸势力；朝廷与官员、边地政权与社群；边地接触区；不把族名或部名当固定疆界。《太宗实录》永乐九年条；《明史·本纪》及相关列传；诏令、奏疏或会典版本的1411（永乐九年）条并未写成宫廷传记；它所确证的是朝廷与官员、边地政权与社群围绕“永乐皇帝派遣一失哈探索北满洲”在账本所能辨认的空间尺度的行动痕迹。 “永乐皇帝派遣一失哈探索北满洲”在1411（永乐九年）之后的人事、诏令或交涉应另随可见文书追索；《太宗实录》永乐九年条；《明史·本纪》及相关列传；诏令、奏疏或会典版本并不能独自裁定动机、密议和受影响者经验。 材料：《太宗实录》永乐九年条；《明史·本纪》及相关列传；诏令、奏疏或会典版本。限度：「永乐皇帝派遣一失哈探索北满洲」可由《太宗实录》永乐九年条；《明史·本纪》及相关列传；诏令、奏疏或会典版本定位到1411（永乐九年），但这些叙述未完整保存边地接触区的执行与受影响者经验，不能把沉默写成事实。
  \item[\eventanchor{ML-1411-004}{1411（永乐九年）：足利义持拒绝永乐皇帝镇压倭寇的请求}] 明与海上、朝贡相关政权；朝廷与官员、军队与卫所；账本可辨空间；地点细节仍须回到所列材料。朝贡名义、边境安全与港口贸易的利益使交涉持续发生。 它为后续册封、互市或海防安排留下文书与跨政权比较线索。 材料：《太宗实录》永乐九年条；《明史·外国传》；《朝鲜王朝实录》、琉球《历代宝案》或海外档案。限度：就「足利义持拒绝永乐皇帝镇压倭寇的请求」而论，《太宗实录》永乐九年条；《明史·外国传》；《朝鲜王朝实录》、琉球《历代宝案》或海外档案提供的是1411（永乐九年）的记录线索；朝廷与官员、军队与卫所在账本可辨空间的实际行动细节不能由这一材料链独自补足。
  \item[\eventanchor{ML-1412-001}{1412（永乐十年）十二月18日：宝藏航行：永乐皇帝下令第四次下西洋}] 明；朝廷与官员；账本可辨空间；地点细节仍须回到所列材料。把“宝藏航行：永乐皇帝下令第四次下西洋”放回1412（永乐十年）十二月18日，《太宗实录》永乐十年条；《明史·本纪》及相关列传；诏令、奏疏或会典版本只让我们看见朝廷与官员、海上与海外政权在账本所能辨认的空间尺度的一段处置链。 读“宝藏航行：永乐皇帝下令第四次下西洋”的余波时，应转向1412（永乐十年）十二月18日之后的任命、批答或地方反应，不能把后出的道德评价倒灌进《太宗实录》永乐十年条；《明史·本纪》及相关列传；诏令、奏疏或会典版本所记的时点。 材料：《太宗实录》永乐十年条；《明史·本纪》及相关列传；诏令、奏疏或会典版本。限度：《太宗实录》永乐十年条；《明史·本纪》及相关列传；诏令、奏疏或会典版本使「宝藏航行：永乐皇帝下令第四次下西洋」的时间位置可追查；至于朝廷与官员在账本可辨空间的命令传递、路线和损失，必须保持待核而不能随编年补写。
  \item[\eventanchor{ML-1412-002}{1412（永乐十年）：明代以炮弹作为弹药。}] 明；朝廷与官员；跨区域行政、资源或军政网络；不假定同步执行。《太宗实录》永乐十年条；《明史·本纪》及相关列传；诏令、奏疏或会典版本的1412（永乐十年）条并未写成宫廷传记；它所确证的是朝廷与官员围绕“明代以炮弹作为弹药”在跨区域的行政、资源与军政网络的行动痕迹。 “明代以炮弹作为弹药”在1412（永乐十年）之后的人事、诏令或交涉应另随可见文书追索；《太宗实录》永乐十年条；《明史·本纪》及相关列传；诏令、奏疏或会典版本并不能独自裁定动机、密议和受影响者经验。 材料：《太宗实录》永乐十年条；《明史·本纪》及相关列传；诏令、奏疏或会典版本。限度：「明代以炮弹作为弹药」虽见于《太宗实录》永乐十年条；《明史·本纪》及相关列传；诏令、奏疏或会典版本，但材料形成者未必观察到跨区域行政、资源或军政网络的全部过程；朝廷与官员的行动范围和受影响群体不能由此单独决定。
  \item[\eventanchor{ML-1412-003}{1412（永乐十年）：北方诸部的册封、互市或出征信息构成本年边疆关系线索，部众名称与军数须保留原材料语境。}] 明；朝廷与官员、军队与卫所；跨区域行政、资源或军政网络；不假定同步执行。朝贡名义、边境安全与港口贸易的利益使交涉持续发生。 它为后续册封、互市或海防安排留下文书与跨政权比较线索。 材料：《太宗实录》永乐十年条；《明史·外国传》；《朝鲜王朝实录》、琉球《历代宝案》或海外档案。限度：桥接条目只标示可回查的年度问题与材料链；实录有中央选择性，崇祯期尤其需要奏疏、地方、朝鲜及满文材料交叉。
  \item[\eventanchor{EM-1413-GUIZHOU-PROVINCE}{1413（永乐十一年）：明置贵州承宣布政使司，调整卫所与土司并存的西南治理}] 明与贵州土司区；朝廷与官员、边地政权与社群；边地接触区；不把族名或部名当固定疆界。黔中交通连接湖广、云南和四川，既有卫所与土司事务需要更稳定的协调层级 贵州成为明朝正式省级行政单位；军屯、驿道和土司互动的成本继续由地方承担 材料：《明太宗实录》永乐十一年条；《明史·地理志》；《贵州通志》建置条与地方碑刻。限度：设省可靠；行政名称不能概括不同土司和山地社群的实际关系
  \item[\eventanchor{EM-1413-ZHENGHE-FOURTH-VOYAGE}{1413（永乐十一年）：郑和第四次下西洋，航程扩展至忽鲁谟斯等地}] 明与印度洋港口政权；朝廷与官员；账本可辨空间；地点细节仍须回到所列材料。东南亚、南亚使节网络已形成，永乐朝试图延伸至印度洋西部重要港口 明使团的外交视野和礼物交换扩大；海外政治秩序仍不能被朝贡表文简化 材料：《明太宗实录》永乐十一年条；马欢《瀛涯胜览》；费信《星槎胜览》。限度：航次和主要目的地可靠；航程细节须区分作者亲历、抄本层次与后出汇编
  \item[\eventanchor{ML-1413-001}{1413（永乐十一年）秋：宝藏航海：郑和从南京出发，按常航路线，新增马尔代夫、比特拉、切特拉岛、霍尔木兹四个航点，并有这样的描述：“四方外国船只，并陆路来往的外国商人，都来到此地，聚集买卖，故此国民皆富”。}] 明；朝廷与官员、地方社会与精英、财政与商业行动者；跨区域行政、资源或军政网络；不假定同步执行。《太宗实录》永乐十一年条；《明史·本纪》及相关列传；诏令、奏疏或会典版本的1413（永乐十一年）秋条并未写成宫廷传记；它所确证的是朝廷与官员、地方社会与精英、财政与商业行动者围绕“宝藏航海：郑和从南京出发，按常航路线，新增马尔代夫、比特拉、切特拉岛、霍尔木兹四个航点，并有这样的描述：“四方外国船只，并陆路来往的外国商人，都来到此地，聚集买卖，故此国民皆富””在跨区域的行政、资源与军政网络的行动痕迹。 “宝藏航海：郑和从南京出发，按常航路线，新增马尔代夫、比特拉、切特拉岛、霍尔木兹四个航点，并有这样的描述：“四方外国船只，并陆路来往的外国商人，都来到此地，聚集买卖，故此国民皆富””在1413（永乐十一年）秋之后的人事、诏令或交涉应另随可见文书追索；《太宗实录》永乐十一年条；《明史·本纪》及相关列传；诏令、奏疏或会典版本并不能独自裁定动机、密议和受影响者经验。 材料：《太宗实录》永乐十一年条；《明史·本纪》及相关列传；诏令、奏疏或会典版本。限度：「宝藏航海：郑和从南京出发，按常航路线，新增马尔代夫、比特拉、切特拉岛、霍尔木兹四个航点，并有这样的描述：“四方外国船只，并陆路来往的外国商人，都来到此地，聚集买卖，故此国民皆富”」可由《太宗实录》永乐十一年条；《明史·本纪》及相关列传；诏令、奏疏或会典版本定位到1413（永乐十一年）秋，但这些叙述未完整保存跨区域行政、资源或军政网络的执行与受影响者经验，不能把沉默写成事实。
  \item[\eventanchor{ML-1413-002}{1413（永乐十一年）：罗本钦波·古实日·罗卓坚赞访问南京}] 明；朝廷与官员；跨区域行政、资源或军政网络；不假定同步执行。“罗本钦波·古实日·罗卓坚赞访问南京”在1413（永乐十一年）进入《太宗实录》永乐十一年条；《明史·本纪》及相关列传；诏令、奏疏或会典版本的叙述，因而可将朝廷与官员与跨区域的行政、资源与军政网络一并放回核验范围。 关于“罗本钦波·古实日·罗卓坚赞访问南京”的实际执行与代价，须从1413（永乐十一年）之后的命令和回应另行检验；一份《太宗实录》永乐十一年条；《明史·本纪》及相关列传；诏令、奏疏或会典版本的编年叙述并不足以代替它们。 材料：《太宗实录》永乐十一年条；《明史·本纪》及相关列传；诏令、奏疏或会典版本。限度：《太宗实录》永乐十一年条；《明史·本纪》及相关列传；诏令、奏疏或会典版本记录「罗本钦波·古实日·罗卓坚赞访问南京」的方式限定了可说范围：可追踪1413（永乐十一年）的行动节点，不能凭此还原跨区域行政、资源或军政网络中朝廷与官员的每一处部署。
  \item[\eventanchor{ML-1413-003}{1413（永乐十一年）：永宁寺碑：明朝派遣伊什哈到努尔干地区军事委员会开设驿站，传播佛教}] 明；朝廷与官员、军队与卫所、宗教与文化行动者；跨区域行政、资源或军政网络；不假定同步执行。《太宗实录》永乐十一年条；《明会典》相关条目；地方志、黄册/契约/河工材料为1413（永乐十一年）的“永宁寺碑：明朝派遣伊什哈到努尔干地区军事委员会开设驿站，传播佛教”划出范围：朝廷与官员、军队与卫所、宗教与文化行动者、跨区域的行政、资源与军政网络以及一项尚待地方材料补足的财政动作。 “永宁寺碑：明朝派遣伊什哈到努尔干地区军事委员会开设驿站，传播佛教”在1413（永乐十一年）形成的是资源调度的核验线，而非完成的财政结论；到达、拖欠与成本仍要借异类材料追踪。 材料：《太宗实录》永乐十一年条；《明会典》相关条目；地方志、黄册/契约/河工材料。限度：《太宗实录》永乐十一年条；《明会典》相关条目；地方志、黄册/契约/河工材料只能把「永宁寺碑：明朝派遣伊什哈到努尔干地区军事委员会开设驿站，传播佛教」置入1413（永乐十一年）的年序；对于跨区域行政、资源或军政网络中朝廷与官员、军队与卫所、宗教与文化行动者的路线、人数和动机，仍须另据地方或对方材料核验。
  \item[\eventanchor{EM-1414-TULA}{1414（永乐十二年）：永乐帝北征瓦剌，在土拉河一带交战}] 明与瓦剌；朝廷与官员、军队与卫所、边地政权与社群；边地接触区；不把族名或部名当固定疆界。瓦剌在草原竞争中增强并影响明廷北方安全；永乐试图阻止其形成优势联盟 瓦剌暂受压力但仍保有机动能力；明廷对草原诸部的分化政策继续依赖军事展示 材料：《明太宗实录》永乐十二年条；《明史·瓦剌传》；17世纪《蒙古源流》瓦剌叙事。限度：亲征和交战可靠；明方所称歼获数与瓦剌内部损失不可等同
  \item[\eventanchor{ML-1414-001}{1414（永乐十二年）三月30日：中国第四次统治越南：陈贵圹被俘}] 明；朝廷与官员；跨区域行政、资源或军政网络；不假定同步执行。朝贡名义、边境安全与港口贸易的利益使交涉持续发生。 它为后续册封、互市或海防安排留下文书与跨政权比较线索。 材料：《太宗实录》永乐十二年条；《明史·外国传》；《朝鲜王朝实录》、琉球《历代宝案》或海外档案。限度：「中国第四次统治越南：陈贵圹被俘」在1414（永乐十二年）三月30日的出现可回查《太宗实录》永乐十二年条；《明史·外国传》；《朝鲜王朝实录》、琉球《历代宝案》或海外档案；它们不等于跨区域行政、资源或军政网络的全景记录，朝廷与官员未被记下的选择与损失不作推定。
  \item[\eventanchor{ML-1414-002}{1414（永乐十二年）四月：第二次蒙古战役：明军在土拉河与卫拉特交战，伤亡惨重，但最终通过重炮轰击取得胜利}] 明与蒙古诸部；朝廷与官员、军队与卫所、边地政权与社群；边地接触区；不把族名或部名当固定疆界。《太宗实录》永乐十二年条；《明会典》相关条目；地方志、黄册/契约/河工材料为1414（永乐十二年）四月的“第二次蒙古战役：明军在土拉河与卫拉特交战，伤亡惨重，但最终通过重炮轰击取得胜利”划出范围：朝廷与官员、军队与卫所、边地政权与社群、账本所能辨认的空间尺度以及一项尚待地方材料补足的财政动作。 “第二次蒙古战役：明军在土拉河与卫拉特交战，伤亡惨重，但最终通过重炮轰击取得胜利”在1414（永乐十二年）四月形成的是资源调度的核验线，而非完成的财政结论；到达、拖欠与成本仍要借异类材料追踪。 材料：《太宗实录》永乐十二年条；《明会典》相关条目；地方志、黄册/契约/河工材料。限度：对「第二次蒙古战役：明军在土拉河与卫拉特交战，伤亡惨重，但最终通过重炮轰击取得胜利」的《太宗实录》永乐十二年条；《明会典》相关条目；地方志、黄册/契约/河工材料记录可支持年序核验；它不足以替代边地接触区的地方材料，故不把朝廷与官员、军队与卫所、边地政权与社群的战报、奏报或叙事当成无争议结果。
  \item[\eventanchor{ML-1414-003}{1414（永乐十二年）：确杰释迦耶喜访问南京}] 明；朝廷与官员；跨区域行政、资源或军政网络；不假定同步执行。《太宗实录》永乐十二年条；《明史·本纪》及相关列传；诏令、奏疏或会典版本为1414（永乐十二年）的“确杰释迦耶喜访问南京”保留了一个窄而具体的观察面：朝廷与官员在跨区域的行政、资源与军政网络的行政动作。 “确杰释迦耶喜访问南京”留给后续政务的是一条从1414（永乐十二年）继续核对的文书次序；不能凭《太宗实录》永乐十二年条；《明史·本纪》及相关列传；诏令、奏疏或会典版本判断谁意图何在或哪一派应负责任。 材料：《太宗实录》永乐十二年条；《明史·本纪》及相关列传；诏令、奏疏或会典版本。限度：以《太宗实录》永乐十二年条；《明史·本纪》及相关列传；诏令、奏疏或会典版本核对「确杰释迦耶喜访问南京」时，只能确认其在1414（永乐十二年）的叙述位置；朝廷与官员在跨区域行政、资源或军政网络的具体经验仍需由可定位的异类材料补证。
  \item[\eventanchor{EM-1415-CANAL-TRANSPORT}{1415（永乐十三年）：会通河航路投入漕粮北运，江南粮食可较稳定转输北京方向}] 明与运河沿线州县；朝廷与官员、地方社会与精英、财政与商业行动者；地方或省域；未具城镇者不补造路线。北京军政中心逐渐成形，海运和陆运难独担粮食供给；河工完成提供新路径 南粮北运成为双都体制的物质基础；河道维护和漕役将长期连接江南财政区与京师 材料：《明太宗实录》永乐十三年漕运条；《明史·食货志》《河渠志》；运河沿线闸坝碑刻。限度：通航与转运政策可靠；实际漕额、损耗和民夫负担须按年度、河段另核
  \item[\eventanchor{ML-1415-001}{1415（永乐十三年）：宝藏航行：宝藏船队捕获了 Sekandar，一名反抗 Samudera Pasai 苏丹国国王 Zain al-'Abidin 的叛军}] 明与地方反抗、流动作战集团；朝廷与官员、军队与卫所、地方社会与精英、流动与反抗集团；地方或省域；未具城镇者不补造路线。对于明与地方反抗、流动作战集团的「宝藏航行：宝藏船队捕获了 Sekandar，一名反抗 Samudera Pasai 苏丹国国王 Zain al-'Abidin 的叛军」，地方或省域与朝廷与官员、军队与卫所、地方社会与精英、流动与反抗集团限定了行动的可行范围；史料未逐项记录前置命令和供给，不能以后来结果补造一个整齐的起因。 「宝藏航行：宝藏船队捕获了 Sekandar，一名反抗 Samudera Pasai 苏丹国国王 Zain al-'Abidin 的叛军」之后，朝廷与官员、军队与卫所、地方社会与精英、流动与反抗集团在地方或省域的行动空间有所改写；这种变化须由随后文书和地方材料追踪，不能把初报数字直接当作最终后果。 材料：《太宗实录》永乐十三年条；《明史·兵志》及相关列传；边镇、地方志或对方材料。限度：关于「宝藏航行：宝藏船队捕获了 Sekandar，一名反抗 Samudera Pasai 苏丹国国王 Zain al-'Abidin 的叛军」，《太宗实录》永乐十三年条；《明史·兵志》及相关列传；边镇、地方志或对方材料所能确证的是条目所示的时点和行动；地方或省域内朝廷与官员、军队与卫所、地方社会与精英、流动与反抗集团的规模、意图及社会后果需要异源材料才能判断。
  \item[\eventanchor{ML-1415-002}{1415（永乐十三年）六月：大运河已重建}] 明；朝廷与官员；跨区域行政、资源或军政网络；不假定同步执行。就1415（永乐十三年）六月而言，“大运河已重建”让朝廷与官员在跨区域的行政、资源与军政网络的资源安排进入《太宗实录》永乐十三年条；《明会典》相关条目；地方志、黄册/契约/河工材料的可见范围。 “大运河已重建”只能把1415（永乐十三年）六月的线索接到地方志、契约、河工或仓储材料；《太宗实录》永乐十三年条；《明会典》相关条目；地方志、黄册/契约/河工材料所载的中央条文不能直接界定执行范围。 材料：《太宗实录》永乐十三年条；《明会典》相关条目；地方志、黄册/契约/河工材料。限度：「大运河已重建」借《太宗实录》永乐十三年条；《明会典》相关条目；地方志、黄册/契约/河工材料获得可检索的时间坐标；跨区域行政、资源或军政网络里朝廷与官员的路线、人数与后果若无其他材料相证，均保留为不可见的部分。
  \item[\eventanchor{ML-1415-003}{1415（永乐十三年）八月12日：宝藏之旅：宝船队返回南京}] 明；朝廷与官员；跨区域行政、资源或军政网络；不假定同步执行。就1415（永乐十三年）八月12日而言，《太宗实录》永乐十三年条；《明史·本纪》及相关列传；诏令、奏疏或会典版本记录“宝藏之旅：宝船队返回南京”时同时划出朝廷与官员与跨区域的行政、资源与军政网络这两个不能省略的条件。 读“宝藏之旅：宝船队返回南京”的余波时，应转向1415（永乐十三年）八月12日之后的任命、批答或地方反应，不能把后出的道德评价倒灌进《太宗实录》永乐十三年条；《明史·本纪》及相关列传；诏令、奏疏或会典版本所记的时点。 材料：《太宗实录》永乐十三年条；《明史·本纪》及相关列传；诏令、奏疏或会典版本。限度：《太宗实录》永乐十三年条；《明史·本纪》及相关列传；诏令、奏疏或会典版本为「宝藏之旅：宝船队返回南京」留下编年入口，却未必保留跨区域行政、资源或军政网络中朝廷与官员的完整视角；本条因此不据之裁定人数、责任或动机。
  \item[\eventanchor{ML-1415-004}{1415（永乐十三年）八月13日：宝藏航行：郑和的同事被派往孟加拉运送礼物}] 明；朝廷与官员；跨区域行政、资源或军政网络；不假定同步执行。“宝藏航行：郑和的同事被派往孟加拉运送礼物”的可证起点不是派系标签，而是1415（永乐十三年）八月13日的《太宗实录》永乐十三年条；《明史·本纪》及相关列传；诏令、奏疏或会典版本对朝廷与官员在跨区域的行政、资源与军政网络所作的登记。 “宝藏航行：郑和的同事被派往孟加拉运送礼物”在1415（永乐十三年）八月13日之后的人事、诏令或交涉应另随可见文书追索；《太宗实录》永乐十三年条；《明史·本纪》及相关列传；诏令、奏疏或会典版本并不能独自裁定动机、密议和受影响者经验。 材料：《太宗实录》永乐十三年条；《明史·本纪》及相关列传；诏令、奏疏或会典版本。限度：《太宗实录》永乐十三年条；《明史·本纪》及相关列传；诏令、奏疏或会典版本记录「宝藏航行：郑和的同事被派往孟加拉运送礼物」的方式限定了可说范围：可追踪1415（永乐十三年）八月13日的行动节点，不能凭此还原跨区域行政、资源或军政网络中朝廷与官员的每一处部署。
  \item[\eventanchor{EM-1416-BEIJING-CONSTRUCTION}{1416（永乐十四年）：北京宫城、坛庙和城防工程持续集中施工}] 明；朝廷与官员、财政与商业行动者；跨区域行政、资源或军政网络；不假定同步执行。迁都需要宫殿、礼制建筑、仓储和驻军空间同时到位；北方木材、砖石和劳力被集中调发 北京具备取代南京为常驻政治中心的物质条件；大规模征发也使工程成本进入地方财政链 材料：《明太宗实录》永乐十四年营建条；《明史·成祖本纪》；北京故宫与城垣考古报告。限度：工程阶段可靠；遗址分期不能精确还原每座建筑的当年完成日
  \item[\eventanchor{ML-1416-001}{1416（永乐十四年）十一月19日：宝藏之旅：永乐皇帝赐礼给18国使节}] 明；朝廷与官员；跨区域行政、资源或军政网络；不假定同步执行。《太宗实录》永乐十四年条；《明史·本纪》及相关列传；诏令、奏疏或会典版本为1416（永乐十四年）十一月19日的“宝藏之旅：永乐皇帝赐礼给18国使节”保留了一个窄而具体的观察面：朝廷与官员在跨区域的行政、资源与军政网络的行政动作。 “宝藏之旅：永乐皇帝赐礼给18国使节”留给后续政务的是一条从1416（永乐十四年）十一月19日继续核对的文书次序；不能凭《太宗实录》永乐十四年条；《明史·本纪》及相关列传；诏令、奏疏或会典版本判断谁意图何在或哪一派应负责任。 材料：《太宗实录》永乐十四年条；《明史·本纪》及相关列传；诏令、奏疏或会典版本。限度：以《太宗实录》永乐十四年条；《明史·本纪》及相关列传；诏令、奏疏或会典版本核对「宝藏之旅：永乐皇帝赐礼给18国使节」时，只能确认其在1416（永乐十四年）十一月19日的叙述位置；朝廷与官员在跨区域行政、资源或军政网络的具体经验仍需由可定位的异类材料补证。
  \item[\eventanchor{ML-1416-002}{1416（永乐十四年）十二月19日：宝藏航行：永乐皇帝下令第五次下西洋}] 明；朝廷与官员；账本可辨空间；地点细节仍须回到所列材料。就1416（永乐十四年）十二月19日而言，《太宗实录》永乐十四年条；《明史·本纪》及相关列传；诏令、奏疏或会典版本记录“宝藏航行：永乐皇帝下令第五次下西洋”时同时划出朝廷与官员、海上与海外政权与账本所能辨认的空间尺度这两个不能省略的条件。 读“宝藏航行：永乐皇帝下令第五次下西洋”的余波时，应转向1416（永乐十四年）十二月19日之后的任命、批答或地方反应，不能把后出的道德评价倒灌进《太宗实录》永乐十四年条；《明史·本纪》及相关列传；诏令、奏疏或会典版本所记的时点。 材料：《太宗实录》永乐十四年条；《明史·本纪》及相关列传；诏令、奏疏或会典版本。限度：《太宗实录》永乐十四年条；《明史·本纪》及相关列传；诏令、奏疏或会典版本只能把「宝藏航行：永乐皇帝下令第五次下西洋」置入1416（永乐十四年）十二月19日的年序；对于账本可辨空间中朝廷与官员的路线、人数和动机，仍须另据地方或对方材料核验。
  \item[\eventanchor{ML-1416-003}{1416（永乐十四年）：迁都筹备、官署调度和礼制安排在本年诏令与奏议中可追踪，文书形成不等于工程即时完成。}] 明；朝廷与官员、财政与商业行动者、宗教与文化行动者；跨区域行政、资源或军政网络；不假定同步执行。“迁都筹备、官署调度和礼制安排在本年诏令与奏议中可追踪，文书形成不等于工程即时完成”的可证起点不是派系标签，而是1416（永乐十四年）的《太宗实录》永乐十四年条；《明史·本纪》及相关列传；诏令、奏疏或会典版本对朝廷与官员、财政与商业行动者、宗教与文化行动者在跨区域的行政、资源与军政网络所作的登记。 “迁都筹备、官署调度和礼制安排在本年诏令与奏议中可追踪，文书形成不等于工程即时完成”在1416（永乐十四年）之后的人事、诏令或交涉应另随可见文书追索；《太宗实录》永乐十四年条；《明史·本纪》及相关列传；诏令、奏疏或会典版本并不能独自裁定动机、密议和受影响者经验。 材料：《太宗实录》永乐十四年条；《明史·本纪》及相关列传；诏令、奏疏或会典版本。限度：桥接条目只标示可回查的年度问题与材料链；实录有中央选择性，崇祯期尤其需要奏疏、地方、朝鲜及满文材料交叉。
  \item[\eventanchor{EM-1417-ZHENGHE-FIFTH-VOYAGE}{1417（永乐十五年）：郑和第五次航行护送使节并继续西抵印度洋港口}] 明与印度洋、东非港口政权；朝廷与官员；账本可辨空间；地点细节仍须回到所列材料。前次航行建立的使节返送需求累积，永乐朝以舰队维持朝贡网络的可见性 海上往来持续扩展明廷的外部信息来源；航行并未改变明朝对民间海外活动的限制框架 材料：《明太宗实录》永乐十五年条；《明史·郑和传》；马欢《瀛涯胜览》。限度：航次可靠；东非来使的具体路线和动物馈赠链需与海外材料互校
  \item[\eventanchor{ML-1417-002}{1417（永乐十五年）秋：宝藏航行：郑和从中国出发，沿先前航线前往霍尔木兹，然后亚丁、摩加迪沙、巴拉瓦、朱布和马林迪}] 明；朝廷与官员；跨区域行政、资源或军政网络；不假定同步执行。“宝藏航行：郑和从中国出发，沿先前航线前往霍尔木兹，然后亚丁、摩加迪沙、巴拉瓦、朱布和马林迪”的可证起点不是派系标签，而是1417（永乐十五年）秋的《太宗实录》永乐十五年条；《明史·本纪》及相关列传；诏令、奏疏或会典版本对朝廷与官员在跨区域的行政、资源与军政网络所作的登记。 “宝藏航行：郑和从中国出发，沿先前航线前往霍尔木兹，然后亚丁、摩加迪沙、巴拉瓦、朱布和马林迪”在1417（永乐十五年）秋之后的人事、诏令或交涉应另随可见文书追索；《太宗实录》永乐十五年条；《明史·本纪》及相关列传；诏令、奏疏或会典版本并不能独自裁定动机、密议和受影响者经验。 材料：《太宗实录》永乐十五年条；《明史·本纪》及相关列传；诏令、奏疏或会典版本。限度：就「宝藏航行：郑和从中国出发，沿先前航线前往霍尔木兹，然后亚丁、摩加迪沙、巴拉瓦、朱布和马林迪」而论，《太宗实录》永乐十五年条；《明史·本纪》及相关列传；诏令、奏疏或会典版本提供的是1417（永乐十五年）秋的记录线索；朝廷与官员在跨区域行政、资源或军政网络的实际行动细节不能由这一材料链独自补足。
  \item[\eventanchor{ML-1417-003}{1417（永乐十五年）：北方诸部的册封、互市或出征信息构成本年边疆关系线索，部众名称与军数须保留原材料语境。}] 明；朝廷与官员、军队与卫所；跨区域行政、资源或军政网络；不假定同步执行。朝贡名义、边境安全与港口贸易的利益使交涉持续发生。 它为后续册封、互市或海防安排留下文书与跨政权比较线索。 材料：《太宗实录》永乐十五年条；《明史·外国传》；《朝鲜王朝实录》、琉球《历代宝案》或海外档案。限度：桥接条目只标示可回查的年度问题与材料链；实录有中央选择性，崇祯期尤其需要奏疏、地方、朝鲜及满文材料交叉。
  \item[\eventanchor{ML-1418-001}{1418（永乐十六年）：靖难后军政接收、北方防务与漕运调发的本年记录为重建政权的连续行动提供编年节点。}] 明；朝廷与官员、军队与卫所、财政与商业行动者；跨区域行政、资源或军政网络；不假定同步执行。围绕「靖难后军政接收、北方防务与漕运调发的本年记录为重建政权的连续行动提供编年节点」，明在跨区域行政、资源或军政网络面对的是朝廷与官员、军队与卫所、财政与商业行动者的现场处置；材料确认行动进入编年，却不足以把此前调兵、补给或地方压力串成确定因果链。 发生「靖难后军政接收、北方防务与漕运调发的本年记录为重建政权的连续行动提供编年节点」后，跨区域行政、资源或军政网络仍是朝廷与官员、军队与卫所、财政与商业行动者必须面对的行动场域；可见结果是问题被带入下一段年序，未见的伤亡、掠夺与居民去留不在本条擅断。 材料：《太宗实录》永乐十六年条；《明史·兵志》及相关列传；边镇、地方志或对方材料。限度：桥接条目只标示可回查的年度问题与材料链；实录有中央选择性，崇祯期尤其需要奏疏、地方、朝鲜及满文材料交叉。
  \item[\eventanchor{ML-1418-002}{1418（永乐十六年）：迁都筹备、官署调度和礼制安排在本年诏令与奏议中可追踪，文书形成不等于工程即时完成。}] 明；朝廷与官员、财政与商业行动者、宗教与文化行动者；跨区域行政、资源或军政网络；不假定同步执行。在1418（永乐十六年），《太宗实录》永乐十六年条；《明史·本纪》及相关列传；诏令、奏疏或会典版本记下“迁都筹备、官署调度和礼制安排在本年诏令与奏议中可追踪，文书形成不等于工程即时完成”；这使朝廷与官员、财政与商业行动者、宗教与文化行动者在跨区域的行政、资源与军政网络的处置留下可核的文书入口。 读“迁都筹备、官署调度和礼制安排在本年诏令与奏议中可追踪，文书形成不等于工程即时完成”的余波时，应转向1418（永乐十六年）之后的任命、批答或地方反应，不能把后出的道德评价倒灌进《太宗实录》永乐十六年条；《明史·本纪》及相关列传；诏令、奏疏或会典版本所记的时点。 材料：《太宗实录》永乐十六年条；《明史·本纪》及相关列传；诏令、奏疏或会典版本。限度：桥接条目只标示可回查的年度问题与材料链；实录有中央选择性，崇祯期尤其需要奏疏、地方、朝鲜及满文材料交叉。
  \item[\eventanchor{ML-1418-003}{1418（永乐十六年）：使团、朝贡港口与航海调度的本年材料记录官方海上秩序，不能以朝贡名义替代实际贸易网络。}] 明；朝廷与官员、财政与商业行动者；账本可辨空间；地点细节仍须回到所列材料。朝贡名义、边境安全与港口贸易的利益使交涉持续发生。 它为后续册封、互市或海防安排留下文书与跨政权比较线索。 材料：《太宗实录》永乐十六年条；《明史·外国传》；《朝鲜王朝实录》、琉球《历代宝案》或海外档案。限度：桥接条目只标示可回查的年度问题与材料链；实录有中央选择性，崇祯期尤其需要奏疏、地方、朝鲜及满文材料交叉。
  \item[\eventanchor{ML-1419-001}{1419（永乐十七年）八月8日：宝藏之旅：宝船队返回中国}] 明；朝廷与官员；跨区域行政、资源或军政网络；不假定同步执行。在1419（永乐十七年）八月8日，《太宗实录》永乐十七年条；《明史·本纪》及相关列传；诏令、奏疏或会典版本记下“宝藏之旅：宝船队返回中国”；这使朝廷与官员在跨区域的行政、资源与军政网络的处置留下可核的文书入口。 读“宝藏之旅：宝船队返回中国”的余波时，应转向1419（永乐十七年）八月8日之后的任命、批答或地方反应，不能把后出的道德评价倒灌进《太宗实录》永乐十七年条；《明史·本纪》及相关列传；诏令、奏疏或会典版本所记的时点。 材料：《太宗实录》永乐十七年条；《明史·本纪》及相关列传；诏令、奏疏或会典版本。限度：「宝藏之旅：宝船队返回中国」虽见于《太宗实录》永乐十七年条；《明史·本纪》及相关列传；诏令、奏疏或会典版本，但材料形成者未必观察到跨区域行政、资源或军政网络的全部过程；朝廷与官员的行动范围和受影响群体不能由此单独决定。
  \item[\eventanchor{ML-1419-002}{1419（永乐十七年）九月20日：宝藏之旅：大使们向明朝宫廷赠送珍奇动物，其中包括孟加拉人从索马里进口的长颈鹿}] 明与蒙古诸部；朝廷与官员、边地政权与社群；边地接触区；不把族名或部名当固定疆界。“宝藏之旅：大使们向明朝宫廷赠送珍奇动物，其中包括孟加拉人从索马里进口的长颈鹿”见于《太宗实录》永乐十七年条；《明史·本纪》及相关列传；诏令、奏疏或会典版本的1419（永乐十七年）九月20日条。它把朝廷与官员、边地政权与社群置于账本所能辨认的空间尺度，而不是抽象的宫廷舞台。 “宝藏之旅：大使们向明朝宫廷赠送珍奇动物，其中包括孟加拉人从索马里进口的长颈鹿”在1419（永乐十七年）九月20日之后的人事、诏令或交涉应另随可见文书追索；《太宗实录》永乐十七年条；《明史·本纪》及相关列传；诏令、奏疏或会典版本并不能独自裁定动机、密议和受影响者经验。 材料：《太宗实录》永乐十七年条；《明史·本纪》及相关列传；诏令、奏疏或会典版本。限度：《太宗实录》永乐十七年条；《明史·本纪》及相关列传；诏令、奏疏或会典版本使「宝藏之旅：大使们向明朝宫廷赠送珍奇动物，其中包括孟加拉人从索马里进口的长颈鹿」的时间位置可追查；至于朝廷与官员、边地政权与社群在边地接触区的命令传递、路线和损失，必须保持待核而不能随编年补写。
  \item[\eventanchor{ML-1419-003}{1419（永乐十七年）十月2日：宝藏航行：下达建造41艘宝藏船的订单}] 明；朝廷与官员；跨区域行政、资源或军政网络；不假定同步执行。《太宗实录》永乐十七年条；《明史·本纪》及相关列传；诏令、奏疏或会典版本把1419（永乐十七年）十月2日的“宝藏航行：下达建造41艘宝藏船的订单”编入年序；能直接追索的是朝廷与官员在跨区域的行政、资源与军政网络所面对的那一项处置。 关于“宝藏航行：下达建造41艘宝藏船的订单”的实际执行与代价，须从1419（永乐十七年）十月2日之后的命令和回应另行检验；一份《太宗实录》永乐十七年条；《明史·本纪》及相关列传；诏令、奏疏或会典版本的编年叙述并不足以代替它们。 材料：《太宗实录》永乐十七年条；《明史·本纪》及相关列传；诏令、奏疏或会典版本。限度：对「宝藏航行：下达建造41艘宝藏船的订单」的《太宗实录》永乐十七年条；《明史·本纪》及相关列传；诏令、奏疏或会典版本记录可支持年序核验；它不足以替代跨区域行政、资源或军政网络的地方材料，故不把朝廷与官员的战报、奏报或叙事当成无争议结果。
  \item[\eventanchor{ML-1419-004}{1419（永乐十七年）：元宵节期间，明皇宫会燃放烟火，火箭沿着电线运行，点燃灯笼，照亮宫殿。}] 明；朝廷与官员；跨区域行政、资源或军政网络；不假定同步执行。沿着1419（永乐十七年）的《太宗实录》永乐十七年条；《明史·本纪》及相关列传；诏令、奏疏或会典版本回看，“元宵节期间，明皇宫会燃放烟火，火箭沿着电线运行，点燃灯笼，照亮宫殿”首先是朝廷与官员在跨区域的行政、资源与军政网络留下的记录，而非后来人物评判的注脚。 “元宵节期间，明皇宫会燃放烟火，火箭沿着电线运行，点燃灯笼，照亮宫殿”令后续核验获得1419（永乐十七年）这一时间坐标；《太宗实录》永乐十七年条；《明史·本纪》及相关列传；诏令、奏疏或会典版本所见的行动者与空间不能推成完整权力因果或统一地方经验。 材料：《太宗实录》永乐十七年条；《明史·本纪》及相关列传；诏令、奏疏或会典版本。限度：《太宗实录》永乐十七年条；《明史·本纪》及相关列传；诏令、奏疏或会典版本记录「元宵节期间，明皇宫会燃放烟火，火箭沿着电线运行，点燃灯笼，照亮宫殿」的方式限定了可说范围：可追踪1419（永乐十七年）的行动节点，不能凭此还原跨区域行政、资源或军政网络中朝廷与官员的每一处部署。
  \item[\eventanchor{EM-1420-BEIJING-CAPITAL}{1420（永乐十八年）：永乐帝正式定北京为京师并完成迁都制度安排}] 明；朝廷与官员；跨区域行政、资源或军政网络；不假定同步执行。北方亲征、北京营建和运河漕运提供了迁都条件；永乐政权需巩固北平政治基础 北京成为皇帝与中枢常驻地，南京保留留都地位；漕运、卫戍和北边决策更紧密相连 材料：《明太宗实录》永乐十八年条；《明史·成祖本纪》《地理志》；北京故宫考古报告。限度：诏令和定都可靠；南京并未失去全部制度功能，不能写作单一中心的即时替换
  \item[\eventanchor{ML-1420-001}{1420（永乐十八年）：紫禁城：天坛建造完成}] 明；朝廷与官员；跨区域行政、资源或军政网络；不假定同步执行。1420（永乐十八年）的“紫禁城：天坛建造完成”由《太宗实录》永乐十八年条；《明史·本纪》及相关列传；诏令、奏疏或会典版本留下线索，记录中的行动者是朝廷与官员，场域是跨区域的行政、资源与军政网络。 “紫禁城：天坛建造完成”留给后续政务的是一条从1420（永乐十八年）继续核对的文书次序；不能凭《太宗实录》永乐十八年条；《明史·本纪》及相关列传；诏令、奏疏或会典版本判断谁意图何在或哪一派应负责任。 材料：《太宗实录》永乐十八年条；《明史·本纪》及相关列传；诏令、奏疏或会典版本。限度：《太宗实录》永乐十八年条；《明史·本纪》及相关列传；诏令、奏疏或会典版本为「紫禁城：天坛建造完成」留下编年入口，却未必保留跨区域行政、资源或军政网络中朝廷与官员的完整视角；本条因此不据之裁定人数、责任或动机。
  \item[\eventanchor{ML-1420-002}{1420（永乐十八年）十月28日：北京正式成为明朝首都}] 明；朝廷与官员；跨区域行政、资源或军政网络；不假定同步执行。把“北京正式成为明朝首都”放回1420（永乐十八年）十月28日，《太宗实录》永乐十八年条；《明史·本纪》及相关列传；诏令、奏疏或会典版本只让我们看见朝廷与官员在跨区域的行政、资源与军政网络的一段处置链。 读“北京正式成为明朝首都”的余波时，应转向1420（永乐十八年）十月28日之后的任命、批答或地方反应，不能把后出的道德评价倒灌进《太宗实录》永乐十八年条；《明史·本纪》及相关列传；诏令、奏疏或会典版本所记的时点。 材料：《太宗实录》永乐十八年条；《明史·本纪》及相关列传；诏令、奏疏或会典版本。限度：《太宗实录》永乐十八年条；《明史·本纪》及相关列传；诏令、奏疏或会典版本记录「北京正式成为明朝首都」的方式限定了可说范围：可追踪1420（永乐十八年）十月28日的行动节点，不能凭此还原跨区域行政、资源或军政网络中朝廷与官员的每一处部署。
  \item[\eventanchor{EM-1421-PALACE-FIRE}{1421（永乐十九年）五月：北京奉天、华盖、谨身三殿遭火灾}] 明；朝廷与官员；跨区域行政、资源或军政网络；不假定同步执行。新都木构宫殿密集且防火条件有限；朝廷以灾异语言讨论君德与政务 中枢礼仪空间受损并需重建；火灾成为反思营建和北征成本的政治话题 材料：《明太宗实录》永乐十九年五月条；《明史·五行志》；北京故宫遗址发掘简报（火灾层）。限度：火灾发生可靠；灾异解释属于政治修辞而非可验证因果
  \item[\eventanchor{ML-1421-001}{1421（永乐十九年）三月3日：宝藏航行：第六次下西洋的命令，向霍尔木兹等16个国家的使节赠送纸币、钱币、礼服和衬里等礼物}] 明；朝廷与官员、财政与商业行动者；账本可辨空间；地点细节仍须回到所列材料。把“宝藏航行：第六次下西洋的命令，向霍尔木兹等16个国家的使节赠送纸币、钱币、礼服和衬里等礼物”还原到1421（永乐十九年）三月3日，《太宗实录》永乐十九年条；《明会典》相关条目；地方志、黄册/契约/河工材料显示朝廷与官员、海上与海外政权、财政与商业行动者在账本所能辨认的空间尺度应对的不是无地点的制度，而是具体资源环节。 从“宝藏航行：第六次下西洋的命令，向霍尔木兹等16个国家的使节赠送纸币、钱币、礼服和衬里等礼物”出发，下一步须将制度规定、行政传达和地方承受分开；《太宗实录》永乐十九年条；《明会典》相关条目；地方志、黄册/契约/河工材料在1421（永乐十九年）三月3日的账面语言不能直接充当社会结果。 材料：《太宗实录》永乐十九年条；《明会典》相关条目；地方志、黄册/契约/河工材料。限度：《太宗实录》永乐十九年条；《明会典》相关条目；地方志、黄册/契约/河工材料使「宝藏航行：第六次下西洋的命令，向霍尔木兹等16个国家的使节赠送纸币、钱币、礼服和衬里等礼物」的时间位置可追查；至于朝廷与官员、财政与商业行动者在账本可辨空间的命令传递、路线和损失，必须保持待核而不能随编年补写。
  \item[\eventanchor{ML-1421-002}{1421（永乐十九年）五月14日：宝藏航行：永乐皇帝下令暂停宝藏航行}] 明；朝廷与官员；跨区域行政、资源或军政网络；不假定同步执行。《太宗实录》永乐十九年条；《明史·本纪》及相关列传；诏令、奏疏或会典版本的1421（永乐十九年）五月14日条并未写成宫廷传记；它所确证的是朝廷与官员围绕“宝藏航行：永乐皇帝下令暂停宝藏航行”在跨区域的行政、资源与军政网络的行动痕迹。 “宝藏航行：永乐皇帝下令暂停宝藏航行”在1421（永乐十九年）五月14日之后的人事、诏令或交涉应另随可见文书追索；《太宗实录》永乐十九年条；《明史·本纪》及相关列传；诏令、奏疏或会典版本并不能独自裁定动机、密议和受影响者经验。 材料：《太宗实录》永乐十九年条；《明史·本纪》及相关列传；诏令、奏疏或会典版本。限度：《太宗实录》永乐十九年条；《明史·本纪》及相关列传；诏令、奏疏或会典版本只能把「宝藏航行：永乐皇帝下令暂停宝藏航行」置入1421（永乐十九年）五月14日的年序；对于跨区域行政、资源或军政网络中朝廷与官员的路线、人数和动机，仍须另据地方或对方材料核验。
  \item[\eventanchor{ML-1421-003}{1421（永乐十九年）十一月10日：宝藏航行：向郑和下令，提供洪堡及十六国使节回国的通道；宝藏船队按常规航线前往锡兰，在那里分头前往马尔代夫、霍尔木兹、阿拉伯国家乔法尔、拉萨和亚丁，以及两个非洲国家摩加迪沙和巴拉瓦；郑和访问甘八里}] 明；朝廷与官员；跨区域行政、资源或军政网络；不假定同步执行。“宝藏航行：向郑和下令，提供洪堡及十六国使节回国的通道；宝藏船队按常规航线前往锡兰，在那里分头前往马尔代夫、霍尔木兹、阿拉伯国家乔法尔、拉萨和亚丁，以及两个非洲国家摩加迪沙和巴拉瓦；郑和访问甘八里”在1421（永乐十九年）十一月10日进入《太宗实录》永乐十九年条；《明史·本纪》及相关列传；诏令、奏疏或会典版本的叙述，因而可将朝廷与官员与跨区域的行政、资源与军政网络一并放回核验范围。 关于“宝藏航行：向郑和下令，提供洪堡及十六国使节回国的通道；宝藏船队按常规航线前往锡兰，在那里分头前往马尔代夫、霍尔木兹、阿拉伯国家乔法尔、拉萨和亚丁，以及两个非洲国家摩加迪沙和巴拉瓦；郑和访问甘八里”的实际执行与代价，须从1421（永乐十九年）十一月10日之后的命令和回应另行检验；一份《太宗实录》永乐十九年条；《明史·本纪》及相关列传；诏令、奏疏或会典版本的编年叙述并不足以代替它们。 材料：《太宗实录》永乐十九年条；《明史·本纪》及相关列传；诏令、奏疏或会典版本。限度：「宝藏航行：向郑和下令，提供洪堡及十六国使节回国的通道；宝藏船队按常规航线前往锡兰，在那里分头前往马尔代夫、霍尔木兹、阿拉伯国家乔法尔、拉萨和亚丁，以及两个非洲国家摩加迪沙和巴拉瓦；郑和访问甘八里」在1421（永乐十九年）十一月10日的出现可回查《太宗实录》永乐十九年条；《明史·本纪》及相关列传；诏令、奏疏或会典版本；它们不等于跨区域行政、资源或军政网络的全景记录，朝廷与官员未被记下的选择与损失不作推定。
  \item[\eventanchor{EM-1422-SECOND-NORTHERN-CAMPAIGN}{1422（永乐二十年）：永乐帝再次北征，寻击阿鲁台而未获决定性会战}] 明与阿鲁台部；朝廷与官员、军队与卫所；跨区域行政、资源或军政网络；不假定同步执行。鞑靼势力重组并袭扰边境；永乐试图以亲征维持威势和盟友关系 明军耗费大量转输而难迫使敌方决战；边防更显示出长期消耗性质 材料：《明太宗实录》永乐二十年条；《明史·鞑靼传》；17世纪《蒙古源流》鞑靼叙事。限度：出征可靠；“无功”或“震慑成功”均是不同政治叙事，需区分
  \item[\eventanchor{ML-1422-001}{1422（永乐二十年）：宝藏航行：宝藏船队在Samudera Pasai Sultanate重新集结并访问暹罗，然后返回中国}] 明；朝廷与官员；账本可辨空间；地点细节仍须回到所列材料。《太宗实录》永乐二十年条；《明史·本纪》及相关列传；诏令、奏疏或会典版本的1422（永乐二十年）条并未写成宫廷传记；它所确证的是朝廷与官员、海上与海外政权围绕“宝藏航行：宝藏船队在Samudera Pasai Sultanate重新集结并访问暹罗，然后返回中国”在账本所能辨认的空间尺度的行动痕迹。 “宝藏航行：宝藏船队在Samudera Pasai Sultanate重新集结并访问暹罗，然后返回中国”在1422（永乐二十年）之后的人事、诏令或交涉应另随可见文书追索；《太宗实录》永乐二十年条；《明史·本纪》及相关列传；诏令、奏疏或会典版本并不能独自裁定动机、密议和受影响者经验。 材料：《太宗实录》永乐二十年条；《明史·本纪》及相关列传；诏令、奏疏或会典版本。限度：关于「宝藏航行：宝藏船队在Samudera Pasai Sultanate重新集结并访问暹罗，然后返回中国」，《太宗实录》永乐二十年条；《明史·本纪》及相关列传；诏令、奏疏或会典版本所能确证的是条目所示的时点和行动；账本可辨空间内朝廷与官员的规模、意图及社会后果需要异源材料才能判断。
  \item[\eventanchor{ML-1422-002}{1422（永乐二十年）四月：第三次蒙古战役：明朝军队出兵攻打阿鲁台，但未能与他交战并返回北京}] 明与蒙古诸部；朝廷与官员、军队与卫所、边地政权与社群；边地接触区；不把族名或部名当固定疆界。「第三次蒙古战役：明朝军队出兵攻打阿鲁台，但未能与他交战并返回北京」出现时，朝廷与官员、军队与卫所、边地政权与社群是在边地接触区的限制下行动；资料所给的是年序定位而非动机自述，前期军政压力只能作为解释线索，不能充作定论。 在边地接触区，朝廷与官员、军队与卫所、边地政权与社群因「第三次蒙古战役：明朝军队出兵攻打阿鲁台，但未能与他交战并返回北京」而要重新安排守备、行军或沟通；这一后续压力可被标示，但战果与伤亡不能只凭一方叙事定量。 材料：《太宗实录》永乐二十年条；《明史·兵志》及相关列传；边镇、地方志或对方材料。限度：「第三次蒙古战役：明朝军队出兵攻打阿鲁台，但未能与他交战并返回北京」借《太宗实录》永乐二十年条；《明史·兵志》及相关列传；边镇、地方志或对方材料获得可检索的时间坐标；边地接触区里朝廷与官员、军队与卫所、边地政权与社群的路线、人数与后果若无其他材料相证，均保留为不可见的部分。
  \item[\eventanchor{EM-1423-THIRD-NORTHERN-CAMPAIGN}{1423（永乐二十一年）：永乐帝第三次北征并继续搜寻阿鲁台部众}] 明与阿鲁台部；朝廷与官员、军队与卫所；跨区域行政、资源或军政网络；不假定同步执行。前次远征未能稳定草原关系；朝廷仍将皇帝出征视为北边威慑手段 持续动员加重军粮和马匹供给；皇帝离京与太子监国的分工更加制度化 材料：《明太宗实录》永乐二十一年条；《明史·成祖本纪》与《鞑靼传》；17世纪《蒙古源流》鞑靼叙事。限度：出征与行程可靠；敌方位置和战果受明方行军报告限制
  \item[\eventanchor{ML-1423-001}{1423（永乐二十一年）八月：第四次蒙古战役：永乐皇帝向阿鲁台发起进攻，却发现他已经被瓦剌击败}] 明与蒙古诸部；朝廷与官员、军队与卫所、边地政权与社群；边地接触区；不把族名或部名当固定疆界。「第四次蒙古战役：永乐皇帝向阿鲁台发起进攻，却发现他已经被瓦剌击败」使朝廷与官员、军队与卫所、边地政权与社群在边地接触区的处境变得可见；本条未保存从征发到出动的全链条，故以可核条件而非概括性“压力”标示其背景。 对明与蒙古诸部而言，「第四次蒙古战役：永乐皇帝向阿鲁台发起进攻，却发现他已经被瓦剌击败」使边地接触区的资源和安全议题进入下一轮处置；本条所能承受的是这一转折，不是对战场损失或地方反应的完整裁断。 材料：《太宗实录》永乐二十一年条；《明史·兵志》及相关列传；边镇、地方志或对方材料。限度：「第四次蒙古战役：永乐皇帝向阿鲁台发起进攻，却发现他已经被瓦剌击败」在1423（永乐二十一年）八月的出现可回查《太宗实录》永乐二十一年条；《明史·兵志》及相关列传；边镇、地方志或对方材料；它们不等于边地接触区的全景记录，朝廷与官员、军队与卫所、边地政权与社群未被记下的选择与损失不作推定。
  \item[\eventanchor{ML-1423-002}{1423（永乐二十一年）：永乐时期典籍、学校与礼仪项目的本年记录提供知识工程线索，编纂、刻印和传播应分别核验。}] 明；朝廷与官员、财政与商业行动者、宗教与文化行动者；跨区域行政、资源或军政网络；不假定同步执行。制度、教育或知识传播的需要推动了相关行动或文本形成。 它提供制度和传播史的锚点，不能据此推断社会整体接受。 材料：《太宗实录》永乐二十一年条；《明史·选举志》或《艺文志》；文集、碑刻或书目材料。限度：桥接条目只标示可回查的年度问题与材料链；实录有中央选择性，崇祯期尤其需要奏疏、地方、朝鲜及满文材料交叉。
  \item[\eventanchor{EM-1424-HONGXI-REVERSALS}{1424（洪熙元年）：仁宗停止部分远征和营建项目，赦免、召还一批永乐朝获罪者}] 明；朝廷与官员、军队与卫所；跨区域行政、资源或军政网络；不假定同步执行。北征、营建和交趾占领形成财政与人事压力；新帝依赖东宫旧臣重整政务 朝廷暂缓高成本项目并提高文官影响；政策急转也受仁宗短暂在位限制 材料：《仁宗实录》洪熙元年条；《明史·仁宗本纪》；杨士奇文集相关奏议。限度：诏令可靠；政策是否落实到各地和获释规模需逐件核查
  \item[\eventanchor{EM-1424-YONGLE-DEATH}{1424（永乐二十二年）七月：永乐帝北征归途中去世，太子朱高炽继位}] 明；朝廷与官员、军队与卫所；跨区域行政、资源或军政网络；不假定同步执行。永乐晚年仍以北征应对草原局势，长距离行军加剧君主健康和继承风险 洪熙朝迅速调整永乐时期的工程和军事政策；太子监国网络转为正式政府 材料：《仁宗实录》洪熙元年追叙；《明史·成祖本纪》《仁宗本纪》；明长陵神道碑与陵寝考古报告。限度：死亡与继位可靠；秘不发丧、随行决策等细节多依后出叙事，须慎用
  \item[\eventanchor{ML-1424-001}{1424（永乐二十二年）二月27日：宝藏之旅：郑和出使巨港，授予史进清之子史继孙“纱帽、金纹锦袍、银印”}] 明、后金（清）与辽东诸势力；朝廷与官员、边地政权与社群、财政与商业行动者；边地接触区；不把族名或部名当固定疆界。《太宗实录》永乐二十二年条；《明会典》相关条目；地方志、黄册/契约/河工材料为1424（永乐二十二年）二月27日的“宝藏之旅：郑和出使巨港，授予史进清之子史继孙“纱帽、金纹锦袍、银印””划出范围：朝廷与官员、边地政权与社群、财政与商业行动者、账本所能辨认的空间尺度以及一项尚待地方材料补足的财政动作。 “宝藏之旅：郑和出使巨港，授予史进清之子史继孙“纱帽、金纹锦袍、银印””在1424（永乐二十二年）二月27日形成的是资源调度的核验线，而非完成的财政结论；到达、拖欠与成本仍要借异类材料追踪。 材料：《太宗实录》永乐二十二年条；《明会典》相关条目；地方志、黄册/契约/河工材料。限度：对「宝藏之旅：郑和出使巨港，授予史进清之子史继孙“纱帽、金纹锦袍、银印”」的《太宗实录》永乐二十二年条；《明会典》相关条目；地方志、黄册/契约/河工材料记录可支持年序核验；它不足以替代边地接触区的地方材料，故不把朝廷与官员、边地政权与社群、财政与商业行动者的战报、奏报或叙事当成无争议结果。
  \item[\eventanchor{ML-1424-002}{1424（永乐二十二年）四月：第五次蒙古战役：永乐皇帝率军远征，讨伐阿鲁台部落的残余势力，但未能找到他们}] 明与蒙古诸部；朝廷与官员、军队与卫所、边地政权与社群；边地接触区；不把族名或部名当固定疆界。「第五次蒙古战役：永乐皇帝率军远征，讨伐阿鲁台部落的残余势力，但未能找到他们」使朝廷与官员、军队与卫所、边地政权与社群在边地接触区的处境变得可见；本条未保存从征发到出动的全链条，故以可核条件而非概括性“压力”标示其背景。 「第五次蒙古战役：永乐皇帝率军远征，讨伐阿鲁台部落的残余势力，但未能找到他们」使明与蒙古诸部在边地接触区的下一步选择重新围绕朝廷与官员、军队与卫所、边地政权与社群展开；其影响应落到随后可见的军令、安置或交涉，而非以未经互校的人数概括。 材料：《太宗实录》永乐二十二年条；《明史·兵志》及相关列传；边镇、地方志或对方材料。限度：就「第五次蒙古战役：永乐皇帝率军远征，讨伐阿鲁台部落的残余势力，但未能找到他们」而论，《太宗实录》永乐二十二年条；《明史·兵志》及相关列传；边镇、地方志或对方材料提供的是1424（永乐二十二年）四月的记录线索；朝廷与官员、军队与卫所、边地政权与社群在边地接触区的实际行动细节不能由这一材料链独自补足。
  \item[\eventanchor{ML-1424-003}{1424（永乐二十二年）八月12日：永乐皇帝驾崩}] 明；朝廷与官员；跨区域行政、资源或军政网络；不假定同步执行。就1424（永乐二十二年）八月12日而言，《太宗实录》永乐二十二年条；《明史·本纪》及相关列传；诏令、奏疏或会典版本记录“永乐皇帝驾崩”时同时划出朝廷与官员与跨区域的行政、资源与军政网络这两个不能省略的条件。 读“永乐皇帝驾崩”的余波时，应转向1424（永乐二十二年）八月12日之后的任命、批答或地方反应，不能把后出的道德评价倒灌进《太宗实录》永乐二十二年条；《明史·本纪》及相关列传；诏令、奏疏或会典版本所记的时点。 材料：《太宗实录》永乐二十二年条；《明史·本纪》及相关列传；诏令、奏疏或会典版本。限度：「永乐皇帝驾崩」虽见于《太宗实录》永乐二十二年条；《明史·本纪》及相关列传；诏令、奏疏或会典版本，但材料形成者未必观察到跨区域行政、资源或军政网络的全部过程；朝廷与官员的行动范围和受影响群体不能由此单独决定。
  \item[\eventanchor{ML-1424-004}{1424（永乐二十二年）九月7日：宝藏之旅：朱高炽即位洪熙皇帝，宝藏之旅结束}] 明；朝廷与官员；跨区域行政、资源或军政网络；不假定同步执行。“宝藏之旅：朱高炽即位洪熙皇帝，宝藏之旅结束”的可证起点不是派系标签，而是1424（永乐二十二年）九月7日的《太宗实录》永乐二十二年条；《明史·本纪》及相关列传；诏令、奏疏或会典版本对朝廷与官员在跨区域的行政、资源与军政网络所作的登记。 “宝藏之旅：朱高炽即位洪熙皇帝，宝藏之旅结束”在1424（永乐二十二年）九月7日之后的人事、诏令或交涉应另随可见文书追索；《太宗实录》永乐二十二年条；《明史·本纪》及相关列传；诏令、奏疏或会典版本并不能独自裁定动机、密议和受影响者经验。 材料：《太宗实录》永乐二十二年条；《明史·本纪》及相关列传；诏令、奏疏或会典版本。限度：「宝藏之旅：朱高炽即位洪熙皇帝，宝藏之旅结束」借《太宗实录》永乐二十二年条；《明史·本纪》及相关列传；诏令、奏疏或会典版本获得可检索的时间坐标；跨区域行政、资源或军政网络里朝廷与官员的路线、人数与后果若无其他材料相证，均保留为不可见的部分。
  \item[\eventanchor{ML-1424-005}{1424（永乐二十二年）：都市毕业生担任县长}] 明；朝廷与官员、地方社会与精英；地方或省域；未具城镇者不补造路线。在1424（永乐二十二年），《太宗实录》永乐二十二年条；《明史·本纪》及相关列传；诏令、奏疏或会典版本记下“都市毕业生担任县长”；这使朝廷与官员、地方社会与精英在所列地方或省域的处置留下可核的文书入口。 读“都市毕业生担任县长”的余波时，应转向1424（永乐二十二年）之后的任命、批答或地方反应，不能把后出的道德评价倒灌进《太宗实录》永乐二十二年条；《明史·本纪》及相关列传；诏令、奏疏或会典版本所记的时点。 材料：《太宗实录》永乐二十二年条；《明史·本纪》及相关列传；诏令、奏疏或会典版本。限度：《太宗实录》永乐二十二年条；《明史·本纪》及相关列传；诏令、奏疏或会典版本为「都市毕业生担任县长」留下编年入口，却未必保留地方或省域中朝廷与官员、地方社会与精英的完整视角；本条因此不据之裁定人数、责任或动机。
\end{description}
